% Options for packages loaded elsewhere
\PassOptionsToPackage{unicode}{hyperref}
\PassOptionsToPackage{hyphens}{url}
\PassOptionsToPackage{dvipsnames,svgnames,x11names}{xcolor}
%
\documentclass[
  10pt,
]{article}
\usepackage{amsmath,amssymb}
\usepackage{iftex}
\ifPDFTeX
  \usepackage[T1]{fontenc}
  \usepackage[utf8]{inputenc}
  \usepackage{textcomp} % provide euro and other symbols
\else % if luatex or xetex
  \usepackage{unicode-math} % this also loads fontspec
  \defaultfontfeatures{Scale=MatchLowercase}
  \defaultfontfeatures[\rmfamily]{Ligatures=TeX,Scale=1}
\fi
\ifPDFTeX\else
  % xetex/luatex font selection
  \setmainfont[]{DejaVu Sans}
  \setmonofont[]{DejaVu Sans Mono}
\fi
% Use upquote if available, for straight quotes in verbatim environments
\IfFileExists{upquote.sty}{\usepackage{upquote}}{}
\IfFileExists{microtype.sty}{% use microtype if available
  \usepackage[]{microtype}
  \UseMicrotypeSet[protrusion]{basicmath} % disable protrusion for tt fonts
}{}
\makeatletter
\@ifundefined{KOMAClassName}{% if non-KOMA class
  \IfFileExists{parskip.sty}{%
    \usepackage{parskip}
  }{% else
    \setlength{\parindent}{0pt}
    \setlength{\parskip}{6pt plus 2pt minus 1pt}}
}{% if KOMA class
  \KOMAoptions{parskip=half}}
\makeatother
\usepackage{xcolor}
\usepackage[margin=2cm]{geometry}
\usepackage{color}
\usepackage{fancyvrb}
\newcommand{\VerbBar}{|}
\newcommand{\VERB}{\Verb[commandchars=\\\{\}]}
\DefineVerbatimEnvironment{Highlighting}{Verbatim}{commandchars=\\\{\}}
% Add ',fontsize=\small' for more characters per line
\newenvironment{Shaded}{}{}
\newcommand{\AlertTok}[1]{\textcolor[rgb]{1.00,0.00,0.00}{\textbf{#1}}}
\newcommand{\AnnotationTok}[1]{\textcolor[rgb]{0.38,0.63,0.69}{\textbf{\textit{#1}}}}
\newcommand{\AttributeTok}[1]{\textcolor[rgb]{0.49,0.56,0.16}{#1}}
\newcommand{\BaseNTok}[1]{\textcolor[rgb]{0.25,0.63,0.44}{#1}}
\newcommand{\BuiltInTok}[1]{\textcolor[rgb]{0.00,0.50,0.00}{#1}}
\newcommand{\CharTok}[1]{\textcolor[rgb]{0.25,0.44,0.63}{#1}}
\newcommand{\CommentTok}[1]{\textcolor[rgb]{0.38,0.63,0.69}{\textit{#1}}}
\newcommand{\CommentVarTok}[1]{\textcolor[rgb]{0.38,0.63,0.69}{\textbf{\textit{#1}}}}
\newcommand{\ConstantTok}[1]{\textcolor[rgb]{0.53,0.00,0.00}{#1}}
\newcommand{\ControlFlowTok}[1]{\textcolor[rgb]{0.00,0.44,0.13}{\textbf{#1}}}
\newcommand{\DataTypeTok}[1]{\textcolor[rgb]{0.56,0.13,0.00}{#1}}
\newcommand{\DecValTok}[1]{\textcolor[rgb]{0.25,0.63,0.44}{#1}}
\newcommand{\DocumentationTok}[1]{\textcolor[rgb]{0.73,0.13,0.13}{\textit{#1}}}
\newcommand{\ErrorTok}[1]{\textcolor[rgb]{1.00,0.00,0.00}{\textbf{#1}}}
\newcommand{\ExtensionTok}[1]{#1}
\newcommand{\FloatTok}[1]{\textcolor[rgb]{0.25,0.63,0.44}{#1}}
\newcommand{\FunctionTok}[1]{\textcolor[rgb]{0.02,0.16,0.49}{#1}}
\newcommand{\ImportTok}[1]{\textcolor[rgb]{0.00,0.50,0.00}{\textbf{#1}}}
\newcommand{\InformationTok}[1]{\textcolor[rgb]{0.38,0.63,0.69}{\textbf{\textit{#1}}}}
\newcommand{\KeywordTok}[1]{\textcolor[rgb]{0.00,0.44,0.13}{\textbf{#1}}}
\newcommand{\NormalTok}[1]{#1}
\newcommand{\OperatorTok}[1]{\textcolor[rgb]{0.40,0.40,0.40}{#1}}
\newcommand{\OtherTok}[1]{\textcolor[rgb]{0.00,0.44,0.13}{#1}}
\newcommand{\PreprocessorTok}[1]{\textcolor[rgb]{0.74,0.48,0.00}{#1}}
\newcommand{\RegionMarkerTok}[1]{#1}
\newcommand{\SpecialCharTok}[1]{\textcolor[rgb]{0.25,0.44,0.63}{#1}}
\newcommand{\SpecialStringTok}[1]{\textcolor[rgb]{0.73,0.40,0.53}{#1}}
\newcommand{\StringTok}[1]{\textcolor[rgb]{0.25,0.44,0.63}{#1}}
\newcommand{\VariableTok}[1]{\textcolor[rgb]{0.10,0.09,0.49}{#1}}
\newcommand{\VerbatimStringTok}[1]{\textcolor[rgb]{0.25,0.44,0.63}{#1}}
\newcommand{\WarningTok}[1]{\textcolor[rgb]{0.38,0.63,0.69}{\textbf{\textit{#1}}}}
\usepackage{longtable,booktabs,array}
\usepackage{calc} % for calculating minipage widths
% Correct order of tables after \paragraph or \subparagraph
\usepackage{etoolbox}
\makeatletter
\patchcmd\longtable{\par}{\if@noskipsec\mbox{}\fi\par}{}{}
\makeatother
% Allow footnotes in longtable head/foot
\IfFileExists{footnotehyper.sty}{\usepackage{footnotehyper}}{\usepackage{footnote}}
\makesavenoteenv{longtable}
\usepackage{graphicx}
\makeatletter
\def\maxwidth{\ifdim\Gin@nat@width>\linewidth\linewidth\else\Gin@nat@width\fi}
\def\maxheight{\ifdim\Gin@nat@height>\textheight\textheight\else\Gin@nat@height\fi}
\makeatother
% Scale images if necessary, so that they will not overflow the page
% margins by default, and it is still possible to overwrite the defaults
% using explicit options in \includegraphics[width, height, ...]{}
\setkeys{Gin}{width=\maxwidth,height=\maxheight,keepaspectratio}
% Set default figure placement to htbp
\makeatletter
\def\fps@figure{htbp}
\makeatother
\setlength{\emergencystretch}{3em} % prevent overfull lines
\providecommand{\tightlist}{%
  \setlength{\itemsep}{0pt}\setlength{\parskip}{0pt}}
\setcounter{secnumdepth}{-\maxdimen} % remove section numbering
\ifLuaTeX
\usepackage[bidi=basic]{babel}
\else
\usepackage[bidi=default]{babel}
\fi
\babelprovide[main,import]{french}
\ifPDFTeX
\else
\babelfont{rm}[]{DejaVu Sans}
\fi
% get rid of language-specific shorthands (see #6817):
\let\LanguageShortHands\languageshorthands
\def\languageshorthands#1{}
\usepackage{fvextra}
\DefineVerbatimEnvironment{Highlighting}{Verbatim}{breaklines,commandchars=\\\{\},fontsize=\scriptsize}
\RecustomVerbatimEnvironment{verbatim}{Verbatim}{breaklines,fontsize=\scriptsize}
\fvset{breaklines=true,fontsize=\scriptsize}
\usepackage{longtable,booktabs,array}
\usepackage{graphicx}
\setlength{\parskip}{4pt}
\AtBeginEnvironment{longtable}{\footnotesize}
\usepackage{etoolbox}
\ifLuaTeX
  \usepackage{selnolig}  % disable illegal ligatures
\fi
\IfFileExists{bookmark.sty}{\usepackage{bookmark}}{\usepackage{hyperref}}
\IfFileExists{xurl.sty}{\usepackage{xurl}}{} % add URL line breaks if available
\urlstyle{same}
\hypersetup{
  pdftitle={F8 --- Peut-on battre la marche aléatoire ?},
  pdfauthor={Groupe : à compléter},
  pdflang={fr},
  colorlinks=true,
  linkcolor={Maroon},
  filecolor={Maroon},
  citecolor={Blue},
  urlcolor={Blue},
  pdfcreator={LaTeX via pandoc}}

\title{F8 --- Peut-on battre la marche aléatoire ?}
\usepackage{etoolbox}
\makeatletter
\providecommand{\subtitle}[1]{% add subtitle to \maketitle
  \apptocmd{\@title}{\par {\large #1 \par}}{}{}
}
\makeatother
\subtitle{Classification supervisée avec référence sur l'indice MASI
(2018-2025)}
\author{Groupe : à compléter}
\date{2026}

\begin{document}
\maketitle

{
\hypersetup{linkcolor=}
\setcounter{tocdepth}{3}
\tableofcontents
}
\begin{longtable}[]{@{}
  >{\raggedright\arraybackslash}p{(\columnwidth - 2\tabcolsep) * \real{0.5000}}
  >{\raggedright\arraybackslash}p{(\columnwidth - 2\tabcolsep) * \real{0.5000}}@{}}
\toprule\noalign{}
\endhead
\bottomrule\noalign{}
\endlastfoot
\textbf{Sujet} & Prévoir le \textbf{sens} (hausse / baisse) du rendement
de la séance suivante de l'indice \textbf{MASI} (Moroccan All Shares
Index) à partir de l'historique de prix. \\
\textbf{Problématique} & Si les prix suivent une marche aléatoire, aucun
modèle ne peut faire mieux qu'une règle triviale. Des algorithmes
d'apprentissage supervisé, nourris d'indicateurs techniques, font-ils
\textbf{significativement} mieux que cette référence, hors échantillon
? \\
\textbf{Objectif mesurable} & Sur le jeu de test chronologique (derniers
15 \% des séances), obtenir une \textbf{précision supérieure à celle de
la référence « marche aléatoire »} (toujours prédire la classe
majoritaire du jeu d'entraînement), avec une p-valeur \textless{} 0,05
(test binomial exact unilatéral). Critères secondaires : précision
équilibrée, MCC, AUC, test de McNemar, validation glissante
(\emph{walk-forward}) et rendement après coûts. \\
\textbf{Données} &
\texttt{data/raw/Moroccan\_All\_Shares\_Historical\_Data\_\_1\_.csv} : 1
988 séances quotidiennes, 02/01/2018 → 31/12/2025. \\
\end{longtable}

\begin{quote}
\textbf{Remarque sur la source.} L'intitulé du sujet mentionne « FRED
SP500 », mais le fichier fourni contient l'indice \textbf{MASI}. Le
notebook est donc réalisé sur le MASI ; seule la cellule de chargement
(section 2) serait à adapter pour le S\&P 500 de FRED (colonne unique
\texttt{SP500}, séparateur décimal différent).
\end{quote}

\textbf{Versions} : affichées dans la cellule de configuration
ci-dessous.

\hypertarget{configuration}{%
\subsection{Configuration}\label{configuration}}

Imports, graine aléatoire (\texttt{GRAINE\ =\ 42}) et chemins relatifs
(le CSV est lu dans \texttt{../data/raw/} ; les figures vont dans
\texttt{../figures/} et la table préparée dans
\texttt{../data/processed/}).

\begin{Shaded}
\begin{Highlighting}[]
\ImportTok{import}\NormalTok{ sys, platform, random, warnings}
\ImportTok{from}\NormalTok{ pathlib }\ImportTok{import}\NormalTok{ Path}
\ImportTok{import}\NormalTok{ numpy }\ImportTok{as}\NormalTok{ np}
\ImportTok{import}\NormalTok{ pandas }\ImportTok{as}\NormalTok{ pd}
\ImportTok{import}\NormalTok{ matplotlib}
\ImportTok{import}\NormalTok{ matplotlib.pyplot }\ImportTok{as}\NormalTok{ plt}
\ImportTok{import}\NormalTok{ scipy}
\ImportTok{from}\NormalTok{ scipy }\ImportTok{import}\NormalTok{ stats}
\ImportTok{import}\NormalTok{ sklearn}
\ImportTok{from}\NormalTok{ sklearn.pipeline }\ImportTok{import}\NormalTok{ Pipeline}
\ImportTok{from}\NormalTok{ sklearn.compose }\ImportTok{import}\NormalTok{ ColumnTransformer}
\ImportTok{from}\NormalTok{ sklearn.preprocessing }\ImportTok{import}\NormalTok{ StandardScaler}
\ImportTok{from}\NormalTok{ sklearn.linear\_model }\ImportTok{import}\NormalTok{ LogisticRegression}
\ImportTok{from}\NormalTok{ sklearn.ensemble }\ImportTok{import}\NormalTok{ RandomForestClassifier, HistGradientBoostingClassifier}
\ImportTok{from}\NormalTok{ sklearn.dummy }\ImportTok{import}\NormalTok{ DummyClassifier}
\ImportTok{from}\NormalTok{ sklearn.metrics }\ImportTok{import}\NormalTok{ (accuracy\_score, balanced\_accuracy\_score, f1\_score,}
\NormalTok{                             matthews\_corrcoef, roc\_auc\_score, confusion\_matrix)}

\NormalTok{warnings.filterwarnings(}\StringTok{"ignore"}\NormalTok{)}
\NormalTok{GRAINE }\OperatorTok{=} \DecValTok{42}
\NormalTok{random.seed(GRAINE)}
\NormalTok{np.random.seed(GRAINE)}

\CommentTok{\# Chemins relatifs : le CSV doit se trouver à côté du notebook (ou dans ./data)}
\NormalTok{NOM\_FICHIER }\OperatorTok{=} \StringTok{"Moroccan\_All\_Shares\_Historical\_Data\_\_1\_.csv"}
\NormalTok{CHEMIN }\OperatorTok{=} \BuiltInTok{next}\NormalTok{(p }\ControlFlowTok{for}\NormalTok{ p }\KeywordTok{in}\NormalTok{ [Path(}\StringTok{"../data/raw"}\NormalTok{) }\OperatorTok{/}\NormalTok{ NOM\_FICHIER, Path(}\StringTok{"data/raw"}\NormalTok{) }\OperatorTok{/}\NormalTok{ NOM\_FICHIER] }\ControlFlowTok{if}\NormalTok{ p.exists())}
\NormalTok{RACINE }\OperatorTok{=}\NormalTok{ CHEMIN.parent.parent.parent                  }\CommentTok{\# racine du projet}
\NormalTok{DOSSIER\_FIG, DOSSIER\_PROC }\OperatorTok{=}\NormalTok{ RACINE }\OperatorTok{/} \StringTok{"figures"}\NormalTok{, RACINE }\OperatorTok{/} \StringTok{"data"} \OperatorTok{/} \StringTok{"processed"}
\NormalTok{DOSSIER\_FIG.mkdir(exist\_ok}\OperatorTok{=}\VariableTok{True}\NormalTok{)}\OperatorTok{;}\NormalTok{ DOSSIER\_PROC.mkdir(exist\_ok}\OperatorTok{=}\VariableTok{True}\NormalTok{)}
\NormalTok{plt.rcParams.update(\{}\StringTok{"figure.figsize"}\NormalTok{: (}\DecValTok{10}\NormalTok{, }\DecValTok{4}\NormalTok{), }\StringTok{"axes.grid"}\NormalTok{: }\VariableTok{True}\NormalTok{, }\StringTok{"grid.alpha"}\NormalTok{: }\FloatTok{0.3}\NormalTok{\})}

\BuiltInTok{print}\NormalTok{(}\StringTok{"Python"}\NormalTok{, platform.python\_version())}
\ControlFlowTok{for}\NormalTok{ nom, mod }\KeywordTok{in}\NormalTok{ [(}\StringTok{"numpy"}\NormalTok{, np), (}\StringTok{"pandas"}\NormalTok{, pd), (}\StringTok{"scipy"}\NormalTok{, scipy), (}\StringTok{"scikit{-}learn"}\NormalTok{, sklearn), (}\StringTok{"matplotlib"}\NormalTok{, matplotlib)]:}
    \BuiltInTok{print}\NormalTok{(}\SpecialStringTok{f"}\SpecialCharTok{\{}\NormalTok{nom}\SpecialCharTok{:13s\}}\SpecialStringTok{"}\NormalTok{, mod.\_\_version\_\_)}
\BuiltInTok{print}\NormalTok{(}\StringTok{"GRAINE ="}\NormalTok{, GRAINE, }\StringTok{"| fichier :"}\NormalTok{, CHEMIN)}
\end{Highlighting}
\end{Shaded}

\begin{Shaded}
\begin{Highlighting}[]
\NormalTok{Python 3.12.3}
\NormalTok{numpy         2.5.4}
\NormalTok{pandas        3.0.6}
\NormalTok{scipy         1.18.1}
\NormalTok{scikit{-}learn  1.9.1}
\NormalTok{matplotlib    3.11.2}
\NormalTok{GRAINE = 42 | fichier : ../data/raw/Moroccan\_All\_Shares\_Historical\_Data\_\_1\_.csv}
\end{Highlighting}
\end{Shaded}

\hypertarget{jeu-de-donnuxe9es}{%
\subsection{2. Jeu de données}\label{jeu-de-donnuxe9es}}

\textbf{Chargement et vérifications} : dimensions, types, premières
lignes, plage de dates et fréquence.

\begin{Shaded}
\begin{Highlighting}[]
\NormalTok{brut }\OperatorTok{=}\NormalTok{ pd.read\_csv(CHEMIN)}
\BuiltInTok{print}\NormalTok{(}\StringTok{"Dimensions brutes :"}\NormalTok{, brut.shape)}
\BuiltInTok{print}\NormalTok{(}\StringTok{"}\CharTok{\textbackslash{}n}\StringTok{Types bruts :"}\NormalTok{)}
\BuiltInTok{print}\NormalTok{(brut.dtypes)}
\NormalTok{display(brut.head())}
\end{Highlighting}
\end{Shaded}

\begin{Shaded}
\begin{Highlighting}[]
\NormalTok{Dimensions brutes : (1988, 7)}

\NormalTok{Types bruts :}
\NormalTok{Date            str}
\NormalTok{Price           str}
\NormalTok{Open            str}
\NormalTok{High            str}
\NormalTok{Low             str}
\NormalTok{Vol.        float64}
\NormalTok{Change \%        str}
\NormalTok{dtype: object}
\NormalTok{         Date      Price       Open       High        Low  Vol. Change \%}
\NormalTok{0  12/31/2025  18,846.35  18,654.74  18,869.80  18,651.55   NaN    1.03\%}
\NormalTok{1  12/30/2025  18,654.74  18,386.49  18,672.90  18,386.49   NaN    1.46\%}
\NormalTok{2  12/29/2025  18,386.49  18,476.19  18,813.23  18,378.58   NaN   {-}0.49\%}
\NormalTok{3  12/26/2025  18,476.19  18,941.16  19,144.14  18,476.19   NaN   {-}2.45\%}
\NormalTok{4  12/25/2025  18,941.16  18,917.37  19,048.87  18,864.61   NaN    0.13\%}
\end{Highlighting}
\end{Shaded}

\begin{Shaded}
\begin{Highlighting}[]
\NormalTok{df }\OperatorTok{=}\NormalTok{ brut.copy()}
\ControlFlowTok{for}\NormalTok{ c }\KeywordTok{in}\NormalTok{ [}\StringTok{"Price"}\NormalTok{, }\StringTok{"Open"}\NormalTok{, }\StringTok{"High"}\NormalTok{, }\StringTok{"Low"}\NormalTok{]:}
\NormalTok{    df[c] }\OperatorTok{=}\NormalTok{ df[c].}\BuiltInTok{str}\NormalTok{.replace(}\StringTok{","}\NormalTok{, }\StringTok{""}\NormalTok{, regex}\OperatorTok{=}\VariableTok{False}\NormalTok{).astype(}\BuiltInTok{float}\NormalTok{)}
\NormalTok{df[}\StringTok{"Change \%"}\NormalTok{] }\OperatorTok{=}\NormalTok{ df[}\StringTok{"Change \%"}\NormalTok{].}\BuiltInTok{str}\NormalTok{.rstrip(}\StringTok{"\%"}\NormalTok{).astype(}\BuiltInTok{float}\NormalTok{)}
\NormalTok{df[}\StringTok{"Date"}\NormalTok{] }\OperatorTok{=}\NormalTok{ pd.to\_datetime(df[}\StringTok{"Date"}\NormalTok{], }\BuiltInTok{format}\OperatorTok{=}\StringTok{"\%m/}\SpecialCharTok{\%d}\StringTok{/\%Y"}\NormalTok{)}
\NormalTok{df }\OperatorTok{=}\NormalTok{ df.sort\_values(}\StringTok{"Date"}\NormalTok{).reset\_index(drop}\OperatorTok{=}\VariableTok{True}\NormalTok{)      }\CommentTok{\# ordre chronologique croissant}

\BuiltInTok{print}\NormalTok{(}\StringTok{"Types après conversion :"}\NormalTok{)}
\BuiltInTok{print}\NormalTok{(df.dtypes)}
\BuiltInTok{print}\NormalTok{(}\StringTok{"}\CharTok{\textbackslash{}n}\StringTok{Plage de dates :"}\NormalTok{, df[}\StringTok{"Date"}\NormalTok{].}\BuiltInTok{min}\NormalTok{().date(), }\StringTok{"→"}\NormalTok{, df[}\StringTok{"Date"}\NormalTok{].}\BuiltInTok{max}\NormalTok{().date())}
\NormalTok{ecarts }\OperatorTok{=}\NormalTok{ df[}\StringTok{"Date"}\NormalTok{].diff().dt.days.dropna()}
\BuiltInTok{print}\NormalTok{(}\StringTok{"Écart médian entre séances :"}\NormalTok{, ecarts.median(), }\StringTok{"jour(s)"}\NormalTok{)}
\BuiltInTok{print}\NormalTok{(}\StringTok{"Répartition des écarts (jours) :"}\NormalTok{)}
\BuiltInTok{print}\NormalTok{(ecarts.value\_counts().sort\_index().to\_string())}
\BuiltInTok{print}\NormalTok{(}\StringTok{"}\CharTok{\textbackslash{}n}\StringTok{Séances par jour de semaine (0 = lundi) :"}\NormalTok{)}
\BuiltInTok{print}\NormalTok{(df[}\StringTok{"Date"}\NormalTok{].dt.dayofweek.value\_counts().sort\_index().to\_string())}
\BuiltInTok{print}\NormalTok{(}\StringTok{"}\CharTok{\textbackslash{}n}\StringTok{Séances par année :"}\NormalTok{)}
\BuiltInTok{print}\NormalTok{(df.groupby(df[}\StringTok{"Date"}\NormalTok{].dt.year).size().to\_string())}
\end{Highlighting}
\end{Shaded}

\begin{Shaded}
\begin{Highlighting}[]
\NormalTok{Types après conversion :}
\NormalTok{Date        datetime64[us]}
\NormalTok{Price              float64}
\NormalTok{Open               float64}
\NormalTok{High               float64}
\NormalTok{Low                float64}
\NormalTok{Vol.               float64}
\NormalTok{Change \%           float64}
\NormalTok{dtype: object}

\NormalTok{Plage de dates : 2018{-}01{-}02 → 2025{-}12{-}31}
\NormalTok{Écart médian entre séances : 1.0 jour(s)}
\NormalTok{Répartition des écarts (jours) :}
\NormalTok{Date}
\NormalTok{1.0    1536}
\NormalTok{2.0      28}
\NormalTok{3.0     383}
\NormalTok{4.0      26}
\NormalTok{5.0      10}
\NormalTok{6.0       3}
\NormalTok{7.0       1}

\NormalTok{Séances par jour de semaine (0 = lundi) :}
\NormalTok{Date}
\NormalTok{0    396}
\NormalTok{1    397}
\NormalTok{2    397}
\NormalTok{3    400}
\NormalTok{4    398}

\NormalTok{Séances par année :}
\NormalTok{Date}
\NormalTok{2018    247}
\NormalTok{2019    246}
\NormalTok{2020    250}
\NormalTok{2021    249}
\NormalTok{2022    254}
\NormalTok{2023    249}
\NormalTok{2024    246}
\NormalTok{2025    247}
\end{Highlighting}
\end{Shaded}

\textbf{Constats.} 1 988 lignes × 7 colonnes ; le CSV est en ordre
antéchronologique (le plus récent en premier) et tous les nombres sont
des chaînes (séparateur de milliers « , », pourcentages avec « \% »).
Après conversion et tri chronologique : période du 02/01/2018 au
31/12/2025, \textbf{fréquence quotidienne en jours de bourse}
(\textasciitilde248 séances par an). Les écarts de 3 jours correspondent
aux week-ends (383 cas), ceux de 4 jours et plus aux jours fériés : il
n'y a pas de séance manquante à reconstituer.

\begin{Shaded}
\begin{Highlighting}[]
\NormalTok{display(df.describe().T.}\BuiltInTok{round}\NormalTok{(}\DecValTok{2}\NormalTok{))}
\BuiltInTok{print}\NormalTok{(}\StringTok{"Valeurs manquantes par colonne :"}\NormalTok{)}
\BuiltInTok{print}\NormalTok{(df.isna().}\BuiltInTok{sum}\NormalTok{().to\_string())}
\end{Highlighting}
\end{Shaded}

\begin{Shaded}
\begin{Highlighting}[]
\NormalTok{           count                        mean                  min                  25\%                  50\%                  75\%                  max          std}
\NormalTok{Date        1988  2021{-}12{-}31 01:10:59.154929  2018{-}01{-}02 00:00:00  2020{-}01{-}07 18:00:00  2022{-}01{-}04 12:00:00  2023{-}12{-}25 06:00:00  2025{-}12{-}31 00:00:00          NaN}
\NormalTok{Price     1988.0                12703.959205              8987.89            11247.465            12107.875           13236.9725              20233.5  2359.926194}
\NormalTok{Open      1988.0                12699.080045              8987.89           11246.9275             12103.69            13232.655              20233.5  2355.860102}
\NormalTok{High      1988.0                12758.151167              9064.97           11296.8825             12160.37            13274.345             20340.23  2378.804814}
\NormalTok{Low       1988.0                12651.332746              8916.12           11199.7525            12057.195            13176.215             20133.01  2343.703562}
\NormalTok{Vol.         0.0                         NaN                  NaN                  NaN                  NaN                  NaN                  NaN          NaN}
\NormalTok{Change \%  1988.0                    0.024074                {-}8.82                 {-}0.3                 0.04                 0.36                 5.45     0.774735}
\NormalTok{Valeurs manquantes par colonne :}
\NormalTok{Date           0}
\NormalTok{Price          0}
\NormalTok{Open           0}
\NormalTok{High           0}
\NormalTok{Low            0}
\NormalTok{Vol.        1988}
\NormalTok{Change \%       0}
\end{Highlighting}
\end{Shaded}

\textbf{Description du jeu (colonnes)}

\begin{longtable}[]{@{}
  >{\raggedright\arraybackslash}p{(\columnwidth - 4\tabcolsep) * \real{0.3333}}
  >{\raggedright\arraybackslash}p{(\columnwidth - 4\tabcolsep) * \real{0.3333}}
  >{\raggedright\arraybackslash}p{(\columnwidth - 4\tabcolsep) * \real{0.3333}}@{}}
\toprule\noalign{}
\begin{minipage}[b]{\linewidth}\raggedright
Colonne
\end{minipage} & \begin{minipage}[b]{\linewidth}\raggedright
Signification
\end{minipage} & \begin{minipage}[b]{\linewidth}\raggedright
Unité / type
\end{minipage} \\
\midrule\noalign{}
\endhead
\bottomrule\noalign{}
\endlastfoot
\texttt{Date} & Date de la séance & jour (format américain
MM/JJ/AAAA) \\
\texttt{Price} & Cours de clôture du MASI & points d'indice \\
\texttt{Open}, \texttt{High}, \texttt{Low} & Ouverture, plus haut, plus
bas & points d'indice \\
\texttt{Vol.} & Volume échangé & \textbf{entièrement vide} (100 \% de
NaN) \\
\texttt{Change\ \%} & Variation en pourcentage de la clôture par rapport
à la veille & \% (arrondi à 2 décimales) \\
\end{longtable}

Le MASI est un indice de capitalisation boursière pondéré par le
flottant, regroupant l'ensemble des actions cotées à la Bourse de
Casablanca. Le fichier est au format des exports d'Investing.com (source
exacte à confirmer par l'auteur). Limites : une seule série, aucun
volume, et une ouverture qui reprend presque toujours la clôture de la
veille.

\hypertarget{exploration}{%
\subsection{3. Exploration}\label{exploration}}

\hypertarget{qualituxe9-des-donnuxe9es}{%
\subsubsection{3.1 Qualité des
données}\label{qualituxe9-des-donnuxe9es}}

Pour chaque problème : méthode de détection → décision → justification.

\begin{Shaded}
\begin{Highlighting}[]
\CommentTok{\# 1) Valeurs manquantes}
\BuiltInTok{print}\NormalTok{(}\StringTok{"Part de NaN dans Vol. :"}\NormalTok{, df[}\StringTok{"Vol."}\NormalTok{].isna().mean())}
\CommentTok{\# 2) Doublons}
\BuiltInTok{print}\NormalTok{(}\StringTok{"Dates dupliquées :"}\NormalTok{, df[}\StringTok{"Date"}\NormalTok{].duplicated().}\BuiltInTok{sum}\NormalTok{(), }\StringTok{"| lignes dupliquées :"}\NormalTok{, df.duplicated().}\BuiltInTok{sum}\NormalTok{())}
\CommentTok{\# 3) Cohérence OHLC}
\NormalTok{inco }\OperatorTok{=}\NormalTok{ ((df[}\StringTok{"High"}\NormalTok{] }\OperatorTok{\textless{}}\NormalTok{ df[}\StringTok{"Low"}\NormalTok{]) }\OperatorTok{|}\NormalTok{ (df[}\StringTok{"Price"}\NormalTok{] }\OperatorTok{\textgreater{}}\NormalTok{ df[}\StringTok{"High"}\NormalTok{]) }\OperatorTok{|}\NormalTok{ (df[}\StringTok{"Price"}\NormalTok{] }\OperatorTok{\textless{}}\NormalTok{ df[}\StringTok{"Low"}\NormalTok{]) }\OperatorTok{|}
\NormalTok{        (df[}\StringTok{"Open"}\NormalTok{] }\OperatorTok{\textgreater{}}\NormalTok{ df[}\StringTok{"High"}\NormalTok{]) }\OperatorTok{|}\NormalTok{ (df[}\StringTok{"Open"}\NormalTok{] }\OperatorTok{\textless{}}\NormalTok{ df[}\StringTok{"Low"}\NormalTok{]))}
\BuiltInTok{print}\NormalTok{(}\StringTok{"Lignes OHLC incohérentes :"}\NormalTok{, }\BuiltInTok{int}\NormalTok{(inco.}\BuiltInTok{sum}\NormalTok{()))}
\CommentTok{\# 4) Cohérence de « Change \% » avec le prix}
\NormalTok{df[}\StringTok{"rendement"}\NormalTok{] }\OperatorTok{=}\NormalTok{ np.log(df[}\StringTok{"Price"}\NormalTok{]).diff()                 }\CommentTok{\# rendement log quotidien}
\NormalTok{ecart\_chg }\OperatorTok{=}\NormalTok{ (df[}\StringTok{"Price"}\NormalTok{].pct\_change() }\OperatorTok{*} \DecValTok{100} \OperatorTok{{-}}\NormalTok{ df[}\StringTok{"Change \%"}\NormalTok{]).}\BuiltInTok{abs}\NormalTok{()}
\BuiltInTok{print}\NormalTok{(}\StringTok{"Écart max |Change \% publié − calculé| :"}\NormalTok{, }\BuiltInTok{round}\NormalTok{(ecart\_chg.}\BuiltInTok{max}\NormalTok{(), }\DecValTok{4}\NormalTok{), }\StringTok{"point de \%  (arrondi à 2 décimales)"}\NormalTok{)}
\CommentTok{\# 5) Open vs clôture précédente}
\NormalTok{gap }\OperatorTok{=}\NormalTok{ (df[}\StringTok{"Open"}\NormalTok{] }\OperatorTok{/}\NormalTok{ df[}\StringTok{"Price"}\NormalTok{].shift(}\DecValTok{1}\NormalTok{) }\OperatorTok{{-}} \DecValTok{1}\NormalTok{) }\OperatorTok{*} \DecValTok{100}
\BuiltInTok{print}\NormalTok{(}\StringTok{"Part de séances où Open == clôture de la veille :"}\NormalTok{, }\BuiltInTok{round}\NormalTok{((gap.}\BuiltInTok{abs}\NormalTok{() }\OperatorTok{\textless{}} \FloatTok{1e{-}9}\NormalTok{).mean(), }\DecValTok{3}\NormalTok{))}
\CommentTok{\# 6) Valeurs atypiques (rendements) : méthode MAD robuste + règle des 5 écarts{-}types}
\NormalTok{r }\OperatorTok{=}\NormalTok{ df[}\StringTok{"rendement"}\NormalTok{].dropna()}
\NormalTok{mad }\OperatorTok{=}\NormalTok{ stats.median\_abs\_deviation(r, scale}\OperatorTok{=}\StringTok{"normal"}\NormalTok{)}
\NormalTok{z\_rob }\OperatorTok{=}\NormalTok{ (r }\OperatorTok{{-}}\NormalTok{ r.median()) }\OperatorTok{/}\NormalTok{ mad}
\NormalTok{masque }\OperatorTok{=}\NormalTok{ z\_rob.}\BuiltInTok{abs}\NormalTok{() }\OperatorTok{\textgreater{}} \DecValTok{5}
\NormalTok{atyp }\OperatorTok{=}\NormalTok{ df.loc[masque[masque].index, [}\StringTok{"Date"}\NormalTok{, }\StringTok{"Price"}\NormalTok{, }\StringTok{"Change \%"}\NormalTok{]].assign(z\_robuste}\OperatorTok{=}\NormalTok{z\_rob[masque].}\BuiltInTok{round}\NormalTok{(}\DecValTok{1}\NormalTok{))}
\BuiltInTok{print}\NormalTok{(}\SpecialStringTok{f"}\CharTok{\textbackslash{}n}\SpecialStringTok{Rendements atypiques (|z robuste| \textgreater{} 5) : }\SpecialCharTok{\{}\BuiltInTok{len}\NormalTok{(atyp)}\SpecialCharTok{\}}\SpecialStringTok{ séances sur }\SpecialCharTok{\{}\BuiltInTok{len}\NormalTok{(r)}\SpecialCharTok{\}}\SpecialStringTok{"}\NormalTok{)}
\BuiltInTok{print}\NormalTok{(atyp.sort\_values(}\StringTok{"Date"}\NormalTok{).to\_string(index}\OperatorTok{=}\VariableTok{False}\NormalTok{))}
\BuiltInTok{print}\NormalTok{(}\StringTok{"Répartition par année :"}\NormalTok{, atyp[}\StringTok{"Date"}\NormalTok{].dt.year.value\_counts().sort\_index().to\_dict())}
\end{Highlighting}
\end{Shaded}

\begin{Shaded}
\begin{Highlighting}[]
\NormalTok{Part de NaN dans Vol. : 1.0}
\NormalTok{Dates dupliquées : 0 | lignes dupliquées : 0}
\NormalTok{Lignes OHLC incohérentes : 0}
\NormalTok{Écart max |Change \% publié − calculé| : 0.0051 point de \%  (arrondi à 2 décimales)}
\NormalTok{Part de séances où Open == clôture de la veille : 0.847}

\NormalTok{Rendements atypiques (|z robuste| \textgreater{} 5) : 31 séances sur 1987}
\NormalTok{      Date    Price  Change \%  z\_robuste}
\NormalTok{2020{-}03{-}06 11474.03     {-}3.05       {-}6.4}
\NormalTok{2020{-}03{-}09 10806.41     {-}5.82      {-}12.3}
\NormalTok{2020{-}03{-}10 11395.21      5.45       10.7}
\NormalTok{2020{-}03{-}12 10579.76     {-}6.70      {-}14.2}
\NormalTok{2020{-}03{-}16  9596.00     {-}8.82      {-}18.9}
\NormalTok{2020{-}03{-}17  9298.54     {-}3.10       {-}6.5}
\NormalTok{2020{-}03{-}18  8987.89     {-}3.34       {-}7.0}
\NormalTok{2020{-}03{-}19  9291.60      3.38        6.7}
\NormalTok{2020{-}03{-}20  9609.27      3.42        6.8}
\NormalTok{2020{-}03{-}23  9953.45      3.58        7.1}
\NormalTok{2020{-}03{-}25  9891.77     {-}2.72       {-}5.7}
\NormalTok{2020{-}06{-}17 10438.48      3.51        6.9}
\NormalTok{2020{-}12{-}21 11164.60     {-}3.09       {-}6.5}
\NormalTok{2022{-}02{-}24 12954.68     {-}4.11       {-}8.6}
\NormalTok{2022{-}03{-}07 12415.46     {-}2.89       {-}6.1}
\NormalTok{2022{-}05{-}10 12666.75     {-}2.44       {-}5.1}
\NormalTok{2023{-}01{-}05 10040.91     {-}3.82       {-}8.0}
\NormalTok{2023{-}01{-}06  9717.99     {-}3.22       {-}6.7}
\NormalTok{2023{-}06{-}20 11704.88      2.63        5.2}
\NormalTok{2023{-}10{-}05 12252.30      5.08       10.0}
\NormalTok{2024{-}08{-}05 13602.42     {-}2.85       {-}6.0}
\NormalTok{2025{-}01{-}03 15516.30      2.84        5.6}
\NormalTok{2025{-}04{-}07 16260.30     {-}5.64      {-}11.9}
\NormalTok{2025{-}04{-}09 16027.71     {-}2.74       {-}5.7}
\NormalTok{2025{-}04{-}14 16579.84      3.17        6.3}
\NormalTok{2025{-}04{-}15 17056.95      2.88        5.7}
\NormalTok{2025{-}06{-}24 18224.19      2.65        5.2}
\NormalTok{2025{-}09{-}09 19521.30     {-}2.59       {-}5.4}
\NormalTok{2025{-}10{-}01 18371.20     {-}3.44       {-}7.2}
\NormalTok{2025{-}10{-}20 18927.90      2.55        5.0}
\NormalTok{2025{-}12{-}26 18476.19     {-}2.45       {-}5.1}
\NormalTok{Répartition par année : \{2020: 13, 2022: 3, 2023: 4, 2024: 1, 2025: 10\}}
\end{Highlighting}
\end{Shaded}

\begin{longtable}[]{@{}
  >{\raggedright\arraybackslash}p{(\columnwidth - 6\tabcolsep) * \real{0.2500}}
  >{\raggedright\arraybackslash}p{(\columnwidth - 6\tabcolsep) * \real{0.2500}}
  >{\raggedright\arraybackslash}p{(\columnwidth - 6\tabcolsep) * \real{0.2500}}
  >{\raggedright\arraybackslash}p{(\columnwidth - 6\tabcolsep) * \real{0.2500}}@{}}
\toprule\noalign{}
\begin{minipage}[b]{\linewidth}\raggedright
Problème
\end{minipage} & \begin{minipage}[b]{\linewidth}\raggedright
Détection
\end{minipage} & \begin{minipage}[b]{\linewidth}\raggedright
Décision
\end{minipage} & \begin{minipage}[b]{\linewidth}\raggedright
Justification
\end{minipage} \\
\midrule\noalign{}
\endhead
\bottomrule\noalign{}
\endlastfoot
\textbf{Valeurs manquantes} & \texttt{isna()} par colonne &
\texttt{Vol.} (100 \% de NaN) \textbf{supprimée} ; aucune autre valeur
manquante & Une colonne vide n'apporte aucune information et ne peut pas
être imputée sans l'inventer. \\
\textbf{Doublons} & Dates et lignes dupliquées & Aucun doublon : rien à
faire & --- \\
\textbf{Incohérences OHLC} & \texttt{Low\ ≤\ Open,\ Price\ ≤\ High} & 0
incohérence & Les prix sont fiables. \\
\textbf{\texttt{Change\ \%} vs prix} & Écart avec la variation
recalculée & Écart max 0,005 point, soit l'arrondi à 2 décimales :
\texttt{Change\ \%} est \textbf{redondante} et n'est pas utilisée ; on
recalcule le rendement log depuis \texttt{Price}. & Éviter une variable
arrondie et mieux contrôlée en interne. \\
\textbf{Ouverture} & Part de séances où \texttt{Open} = clôture de la
veille & 84,7 \% : \texttt{Open} est peu informative ; on ne garde que
l'écart d'ouverture (\emph{gap}) comme variable. & Limite de la source,
à discuter en section 8. \\
\textbf{Valeurs atypiques} & z-score \textbf{robuste} (médiane/MAD) sur
les rendements, seuil \textbar z\textbar{} \textgreater{} 5 & 31 séances
(1,6 \%) signalées, \textbf{toutes conservées} : 13 en 2020 (krach
Covid-19 de mars 2020), 10 en 2025 (dont le choc du 7 avril 2025), etc.
& En finance, une valeur extrême est le plus souvent réelle (krach,
crise). Les supprimer retirerait précisément les épisodes où le risque
se matérialise et biaiserait l'évaluation. On choisit des
modèles/pipelines peu sensibles aux extrêmes (arbres) ou régularisés
(L2). \\
\end{longtable}

\hypertarget{statistiques-descriptives-utiles-au-probluxe8me}{%
\subsubsection{3.2 Statistiques descriptives utiles au
problème}\label{statistiques-descriptives-utiles-au-probluxe8me}}

On étudie les propriétés des \textbf{rendements} : leur équilibre
haussier/baissier, leur éventuelle autocorrélation (si elle est nulle,
la marche aléatoire tient) et leur volatilité.

\textbf{Constats.} (i) 52,8 \% des rendements sont positifs : un
prédicteur « toujours hausse » est déjà un adversaire non trivial ; (ii)
la distribution est très éloignée de la loi normale (asymétrie −1,55,
kurtosis en excès ≈ 20,6, Jarque-Bera rejette la normalité) ; (iii) les
rendements sont \textbf{faiblement mais significativement autocorrélés}
(ρ₁ ≈ 0,126, ρ₂ ≈ 0,114, Ljung-Box à 10 retards : p ≈ 5·10⁻¹¹), ce qui
contredit la marche aléatoire \emph{stricte} ; (iv) les rendements au
carré sont autocorrélés (ρ ≈ 0,20) : la volatilité est regroupée, avec
un pic en mars 2020.

\begin{Shaded}
\begin{Highlighting}[]
\BuiltInTok{print}\NormalTok{(}\StringTok{"Statistiques des rendements log quotidiens :"}\NormalTok{)}
\NormalTok{res }\OperatorTok{=}\NormalTok{ pd.Series(\{}
    \StringTok{"moyenne (\%)"}\NormalTok{: r.mean() }\OperatorTok{*} \DecValTok{100}\NormalTok{, }\StringTok{"écart{-}type (\%)"}\NormalTok{: r.std() }\OperatorTok{*} \DecValTok{100}\NormalTok{,}
    \StringTok{"min (\%)"}\NormalTok{: r.}\BuiltInTok{min}\NormalTok{() }\OperatorTok{*} \DecValTok{100}\NormalTok{, }\StringTok{"max (\%)"}\NormalTok{: r.}\BuiltInTok{max}\NormalTok{() }\OperatorTok{*} \DecValTok{100}\NormalTok{,}
    \StringTok{"asymétrie"}\NormalTok{: stats.skew(r), }\StringTok{"kurtosis excès"}\NormalTok{: stats.kurtosis(r),}
    \StringTok{"part de jours \textgreater{} 0"}\NormalTok{: (r }\OperatorTok{\textgreater{}} \DecValTok{0}\NormalTok{).mean(), }\StringTok{"part de jours = 0"}\NormalTok{: (r }\OperatorTok{==} \DecValTok{0}\NormalTok{).mean()\}).}\BuiltInTok{round}\NormalTok{(}\DecValTok{4}\NormalTok{)}
\BuiltInTok{print}\NormalTok{(res.to\_string())}
\NormalTok{jb }\OperatorTok{=}\NormalTok{ stats.jarque\_bera(r)}
\BuiltInTok{print}\NormalTok{(}\SpecialStringTok{f"}\CharTok{\textbackslash{}n}\SpecialStringTok{Jarque{-}Bera : stat = }\SpecialCharTok{\{}\NormalTok{jb}\SpecialCharTok{.}\NormalTok{statistic}\SpecialCharTok{:.1f\}}\SpecialStringTok{, p{-}valeur = }\SpecialCharTok{\{}\NormalTok{jb}\SpecialCharTok{.}\NormalTok{pvalue}\SpecialCharTok{:.2e\}}\SpecialStringTok{  → normalité rejetée"}\NormalTok{)}
\KeywordTok{def}\NormalTok{ acf(x, nlags):}
\NormalTok{    x }\OperatorTok{=}\NormalTok{ np.asarray(x) }\OperatorTok{{-}}\NormalTok{ np.mean(x)}
    \ControlFlowTok{return}\NormalTok{ np.array([np.}\BuiltInTok{sum}\NormalTok{(x[k:] }\OperatorTok{*}\NormalTok{ x[:}\BuiltInTok{len}\NormalTok{(x) }\OperatorTok{{-}}\NormalTok{ k]) }\OperatorTok{/}\NormalTok{ np.}\BuiltInTok{sum}\NormalTok{(x }\OperatorTok{*}\NormalTok{ x) }\ControlFlowTok{for}\NormalTok{ k }\KeywordTok{in} \BuiltInTok{range}\NormalTok{(}\DecValTok{1}\NormalTok{, nlags }\OperatorTok{+} \DecValTok{1}\NormalTok{)])}
\NormalTok{rho }\OperatorTok{=}\NormalTok{ acf(r, }\DecValTok{10}\NormalTok{)}
\NormalTok{n }\OperatorTok{=} \BuiltInTok{len}\NormalTok{(r)}
\NormalTok{Q }\OperatorTok{=}\NormalTok{ n }\OperatorTok{*}\NormalTok{ (n }\OperatorTok{+} \DecValTok{2}\NormalTok{) }\OperatorTok{*}\NormalTok{ np.}\BuiltInTok{sum}\NormalTok{(rho }\OperatorTok{**} \DecValTok{2} \OperatorTok{/}\NormalTok{ (n }\OperatorTok{{-}}\NormalTok{ np.arange(}\DecValTok{1}\NormalTok{, }\DecValTok{11}\NormalTok{)))}
\BuiltInTok{print}\NormalTok{(}\StringTok{"Autocorrélations lags 1..10 :"}\NormalTok{, np.}\BuiltInTok{round}\NormalTok{(rho, }\DecValTok{3}\NormalTok{))}
\BuiltInTok{print}\NormalTok{(}\SpecialStringTok{f"Ljung{-}Box (10 lags) : Q = }\SpecialCharTok{\{}\NormalTok{Q}\SpecialCharTok{:.1f\}}\SpecialStringTok{, p{-}valeur = }\SpecialCharTok{\{}\DecValTok{1} \OperatorTok{{-}}\NormalTok{ stats}\SpecialCharTok{.}\NormalTok{chi2}\SpecialCharTok{.}\NormalTok{cdf(Q, }\DecValTok{10}\NormalTok{)}\SpecialCharTok{:.2e\}}\SpecialStringTok{"}\NormalTok{)}
\BuiltInTok{print}\NormalTok{(}\SpecialStringTok{f"Autocorrélation des rendements² (lag 1) : }\SpecialCharTok{\{}\NormalTok{acf(r}\OperatorTok{**}\DecValTok{2}\NormalTok{, }\DecValTok{1}\NormalTok{)[}\DecValTok{0}\NormalTok{]}\SpecialCharTok{:.3f\}}\SpecialStringTok{  (regroupement de volatilité)"}\NormalTok{)}
\end{Highlighting}
\end{Shaded}

\begin{Shaded}
\begin{Highlighting}[]
\NormalTok{Statistiques des rendements log quotidiens :}
\NormalTok{moyenne (\%)           0.0210}
\NormalTok{écart{-}type (\%)        0.7791}
\NormalTok{min (\%)              {-}9.2317}
\NormalTok{max (\%)               5.3054}
\NormalTok{asymétrie            {-}1.5497}
\NormalTok{kurtosis excès       20.6384}
\NormalTok{part de jours \textgreater{} 0     0.5279}
\NormalTok{part de jours = 0     0.0005}

\NormalTok{Jarque{-}Bera : stat = 36059.9, p{-}valeur = 0.00e+00  → normalité rejetée}
\NormalTok{Autocorrélations lags 1..10 : [ 0.126  0.114 {-}0.004 {-}0.038 {-}0.022 {-}0.001  0.037 {-}0.049  0.022  0.004]}
\NormalTok{Ljung{-}Box (10 lags) : Q = 69.9, p{-}valeur = 4.58e{-}11}
\NormalTok{Autocorrélation des rendements² (lag 1) : 0.197  (regroupement de volatilité)}
\end{Highlighting}
\end{Shaded}

\hypertarget{graphiques}{%
\subsubsection{3.3 Graphiques}\label{graphiques}}

\emph{Haut gauche} : niveau du MASI (échelle log) et krach de mars 2020.
\emph{Haut droit} : queues épaisses par rapport à la loi normale.
\emph{Bas gauche} : autocorrélations aux retards 1 et 2 au-delà des
bornes de bruit blanc. \emph{Bas droit} : volatilité glissante,
nettement plus forte en 2020.

\begin{Shaded}
\begin{Highlighting}[]
\NormalTok{fig, ax }\OperatorTok{=}\NormalTok{ plt.subplots(}\DecValTok{2}\NormalTok{, }\DecValTok{2}\NormalTok{, figsize}\OperatorTok{=}\NormalTok{(}\DecValTok{13}\NormalTok{, }\DecValTok{8}\NormalTok{))}
\NormalTok{ax[}\DecValTok{0}\NormalTok{, }\DecValTok{0}\NormalTok{].plot(df[}\StringTok{"Date"}\NormalTok{], df[}\StringTok{"Price"}\NormalTok{], lw}\OperatorTok{=}\DecValTok{1}\NormalTok{)}
\NormalTok{ax[}\DecValTok{0}\NormalTok{, }\DecValTok{0}\NormalTok{].set\_yscale(}\StringTok{"log"}\NormalTok{)}
\NormalTok{ax[}\DecValTok{0}\NormalTok{, }\DecValTok{0}\NormalTok{].axvspan(pd.Timestamp(}\StringTok{"2020{-}03{-}01"}\NormalTok{), pd.Timestamp(}\StringTok{"2020{-}04{-}15"}\NormalTok{), color}\OperatorTok{=}\StringTok{"red"}\NormalTok{, alpha}\OperatorTok{=}\FloatTok{.2}\NormalTok{, label}\OperatorTok{=}\StringTok{"krach Covid{-}19"}\NormalTok{)}
\NormalTok{ax[}\DecValTok{0}\NormalTok{, }\DecValTok{0}\NormalTok{].}\BuiltInTok{set}\NormalTok{(title}\OperatorTok{=}\StringTok{"MASI – niveau de clôture (échelle log)"}\NormalTok{, xlabel}\OperatorTok{=}\StringTok{"Date"}\NormalTok{, ylabel}\OperatorTok{=}\StringTok{"Indice (points, échelle log)"}\NormalTok{)}
\NormalTok{ax[}\DecValTok{0}\NormalTok{, }\DecValTok{0}\NormalTok{].legend()}
\NormalTok{ax[}\DecValTok{0}\NormalTok{, }\DecValTok{1}\NormalTok{].hist(r }\OperatorTok{*} \DecValTok{100}\NormalTok{, bins}\OperatorTok{=}\DecValTok{80}\NormalTok{, density}\OperatorTok{=}\VariableTok{True}\NormalTok{, alpha}\OperatorTok{=}\FloatTok{.7}\NormalTok{, label}\OperatorTok{=}\StringTok{"observé"}\NormalTok{)}
\NormalTok{xs }\OperatorTok{=}\NormalTok{ np.linspace(r.}\BuiltInTok{min}\NormalTok{() }\OperatorTok{*} \DecValTok{100}\NormalTok{, r.}\BuiltInTok{max}\NormalTok{() }\OperatorTok{*} \DecValTok{100}\NormalTok{, }\DecValTok{300}\NormalTok{)}
\NormalTok{ax[}\DecValTok{0}\NormalTok{, }\DecValTok{1}\NormalTok{].plot(xs, stats.norm.pdf(xs, r.mean() }\OperatorTok{*} \DecValTok{100}\NormalTok{, r.std() }\OperatorTok{*} \DecValTok{100}\NormalTok{), }\StringTok{"r"}\NormalTok{, label}\OperatorTok{=}\StringTok{"loi normale ajustée"}\NormalTok{)}
\NormalTok{ax[}\DecValTok{0}\NormalTok{, }\DecValTok{1}\NormalTok{].}\BuiltInTok{set}\NormalTok{(title}\OperatorTok{=}\StringTok{"Distribution des rendements log quotidiens"}\NormalTok{, xlabel}\OperatorTok{=}\StringTok{"Rendement (\%)"}\NormalTok{, ylabel}\OperatorTok{=}\StringTok{"Densité"}\NormalTok{)}
\NormalTok{ax[}\DecValTok{0}\NormalTok{, }\DecValTok{1}\NormalTok{].legend()}
\NormalTok{ax[}\DecValTok{1}\NormalTok{, }\DecValTok{0}\NormalTok{].bar(}\BuiltInTok{range}\NormalTok{(}\DecValTok{1}\NormalTok{, }\DecValTok{11}\NormalTok{), rho)}
\NormalTok{borne }\OperatorTok{=} \FloatTok{1.96} \OperatorTok{/}\NormalTok{ np.sqrt(n)}
\NormalTok{ax[}\DecValTok{1}\NormalTok{, }\DecValTok{0}\NormalTok{].axhline(borne, color}\OperatorTok{=}\StringTok{"r"}\NormalTok{, ls}\OperatorTok{=}\StringTok{"{-}{-}"}\NormalTok{)}\OperatorTok{;}\NormalTok{ ax[}\DecValTok{1}\NormalTok{, }\DecValTok{0}\NormalTok{].axhline(}\OperatorTok{{-}}\NormalTok{borne, color}\OperatorTok{=}\StringTok{"r"}\NormalTok{, ls}\OperatorTok{=}\StringTok{"{-}{-}"}\NormalTok{, label}\OperatorTok{=}\StringTok{"bornes 95 \% (bruit blanc)"}\NormalTok{)}
\NormalTok{ax[}\DecValTok{1}\NormalTok{, }\DecValTok{0}\NormalTok{].}\BuiltInTok{set}\NormalTok{(title}\OperatorTok{=}\StringTok{"Autocorrélation des rendements"}\NormalTok{, xlabel}\OperatorTok{=}\StringTok{"Retard (jours de bourse)"}\NormalTok{, ylabel}\OperatorTok{=}\StringTok{"Coefficient d\textquotesingle{}autocorrélation"}\NormalTok{)}
\NormalTok{ax[}\DecValTok{1}\NormalTok{, }\DecValTok{0}\NormalTok{].legend()}
\NormalTok{ax[}\DecValTok{1}\NormalTok{, }\DecValTok{1}\NormalTok{].plot(df[}\StringTok{"Date"}\NormalTok{], df[}\StringTok{"rendement"}\NormalTok{].rolling(}\DecValTok{60}\NormalTok{).std() }\OperatorTok{*}\NormalTok{ np.sqrt(}\DecValTok{252}\NormalTok{) }\OperatorTok{*} \DecValTok{100}\NormalTok{)}
\NormalTok{ax[}\DecValTok{1}\NormalTok{, }\DecValTok{1}\NormalTok{].}\BuiltInTok{set}\NormalTok{(title}\OperatorTok{=}\StringTok{"Volatilité annualisée glissante (60 séances)"}\NormalTok{, xlabel}\OperatorTok{=}\StringTok{"Date"}\NormalTok{, ylabel}\OperatorTok{=}\StringTok{"Volatilité annualisée (\%)"}\NormalTok{)}
\NormalTok{plt.tight\_layout()}\OperatorTok{;}\NormalTok{ plt.savefig(DOSSIER\_FIG }\OperatorTok{/} \StringTok{"01\_exploration.png"}\NormalTok{, dpi}\OperatorTok{=}\DecValTok{150}\NormalTok{, bbox\_inches}\OperatorTok{=}\StringTok{"tight"}\NormalTok{)}\OperatorTok{;}\NormalTok{ plt.show()}
\end{Highlighting}
\end{Shaded}

\begin{figure}
\centering
\includegraphics[width=1\textwidth,height=\textheight]{../../figures/01_exploration.png}
\caption{Exploration du MASI : niveau, distribution, autocorrélation et
volatilité}
\end{figure}

\hypertarget{pruxe9paration}{%
\subsection{4. Préparation}\label{pruxe9paration}}

\begin{itemize}
\tightlist
\item
  \textbf{Nettoyage} : tri chronologique, conversion des types,
  suppression de \texttt{Vol.} et de \texttt{Change\ \%} (redondante) ;
  aucune valeur extrême supprimée (voir 3.1).
\item
  \textbf{Cible} : \texttt{cible\ =\ 1} si le rendement log de la séance
  \textbf{suivante} est strictement positif, 0 sinon (le cas « = 0 » est
  quasi inexistant, 0,05 \%).
\item
  \textbf{Création de variables} (toutes calculables à la clôture de la
  séance \emph{t}, donc sans fuite d'information du futur) : rendements
  retardés \texttt{r\_lag0…r\_lag4}, moyennes mobiles 5 et 20 jours,
  volatilités 5 et 20 jours, amplitude (High−Low)/Close, écart
  d'ouverture, momentum 10 jours, écart à la moyenne mobile 50 jours,
  RSI 14 jours. 50 lignes sont perdues au démarrage (fenêtre de 50
  jours).
\item
  \textbf{Encodage} : le jour de la semaine est encodé en \emph{one-hot}
  (5 indicatrices), car c'est une variable catégorielle sans ordre
  naturel.
\item
  \textbf{Mise à l'échelle} : \textbf{nécessaire pour la régression
  logistique} (la pénalité L2 et l'optimisation dépendent des échelles :
  sans standardisation, une variable en « points » écraserait une
  variable en « \% ») ; \textbf{inutile pour les arbres} (forêt
  aléatoire, gradient boosting), qui ne dépendent que de l'ordre des
  valeurs. Le \texttt{StandardScaler} est \textbf{appris sur le train
  uniquement} (dans un \texttt{Pipeline}) puis appliqué à
  validation/test, pour éviter toute fuite.
\item
  \textbf{Découpage chronologique} 70 \% / 15 \% / 15 \% (pas de mélange
  aléatoire : on prédit le futur avec le passé), avec un \textbf{embargo
  d'une ligne} entre blocs, car la cible de la dernière ligne d'un bloc
  utilise le rendement de la première séance du bloc suivant.
\end{itemize}

\begin{Shaded}
\begin{Highlighting}[]
\NormalTok{d }\OperatorTok{=}\NormalTok{ df.copy()}
\NormalTok{r\_ }\OperatorTok{=}\NormalTok{ d[}\StringTok{"rendement"}\NormalTok{]}
\NormalTok{X }\OperatorTok{=}\NormalTok{ pd.DataFrame(index}\OperatorTok{=}\NormalTok{d.index)}
\ControlFlowTok{for}\NormalTok{ k }\KeywordTok{in} \BuiltInTok{range}\NormalTok{(}\DecValTok{5}\NormalTok{):}
\NormalTok{    X[}\SpecialStringTok{f"r\_lag}\SpecialCharTok{\{}\NormalTok{k}\SpecialCharTok{\}}\SpecialStringTok{"}\NormalTok{] }\OperatorTok{=}\NormalTok{ r\_.shift(k)                     }\CommentTok{\# r\_lag0 = rendement du jour t (connu à la clôture)}
\NormalTok{X[}\StringTok{"moy\_5"}\NormalTok{] }\OperatorTok{=}\NormalTok{ r\_.rolling(}\DecValTok{5}\NormalTok{).mean()}
\NormalTok{X[}\StringTok{"moy\_20"}\NormalTok{] }\OperatorTok{=}\NormalTok{ r\_.rolling(}\DecValTok{20}\NormalTok{).mean()}
\NormalTok{X[}\StringTok{"vol\_5"}\NormalTok{] }\OperatorTok{=}\NormalTok{ r\_.rolling(}\DecValTok{5}\NormalTok{).std()}
\NormalTok{X[}\StringTok{"vol\_20"}\NormalTok{] }\OperatorTok{=}\NormalTok{ r\_.rolling(}\DecValTok{20}\NormalTok{).std()}
\NormalTok{X[}\StringTok{"amplitude"}\NormalTok{] }\OperatorTok{=}\NormalTok{ (d[}\StringTok{"High"}\NormalTok{] }\OperatorTok{{-}}\NormalTok{ d[}\StringTok{"Low"}\NormalTok{]) }\OperatorTok{/}\NormalTok{ d[}\StringTok{"Price"}\NormalTok{]}
\NormalTok{X[}\StringTok{"ecart\_ouverture"}\NormalTok{] }\OperatorTok{=}\NormalTok{ d[}\StringTok{"Open"}\NormalTok{] }\OperatorTok{/}\NormalTok{ d[}\StringTok{"Price"}\NormalTok{].shift(}\DecValTok{1}\NormalTok{) }\OperatorTok{{-}} \DecValTok{1}
\NormalTok{X[}\StringTok{"momentum\_10"}\NormalTok{] }\OperatorTok{=}\NormalTok{ np.log(d[}\StringTok{"Price"}\NormalTok{] }\OperatorTok{/}\NormalTok{ d[}\StringTok{"Price"}\NormalTok{].shift(}\DecValTok{10}\NormalTok{))}
\NormalTok{X[}\StringTok{"ecart\_mm50"}\NormalTok{] }\OperatorTok{=}\NormalTok{ d[}\StringTok{"Price"}\NormalTok{] }\OperatorTok{/}\NormalTok{ d[}\StringTok{"Price"}\NormalTok{].rolling(}\DecValTok{50}\NormalTok{).mean() }\OperatorTok{{-}} \DecValTok{1}
\NormalTok{delta }\OperatorTok{=}\NormalTok{ d[}\StringTok{"Price"}\NormalTok{].diff()}
\NormalTok{gain }\OperatorTok{=}\NormalTok{ delta.clip(lower}\OperatorTok{=}\DecValTok{0}\NormalTok{).rolling(}\DecValTok{14}\NormalTok{).mean()}\OperatorTok{;}\NormalTok{ perte }\OperatorTok{=}\NormalTok{ (}\OperatorTok{{-}}\NormalTok{delta.clip(upper}\OperatorTok{=}\DecValTok{0}\NormalTok{)).rolling(}\DecValTok{14}\NormalTok{).mean()}
\NormalTok{X[}\StringTok{"rsi\_14"}\NormalTok{] }\OperatorTok{=} \DecValTok{100} \OperatorTok{{-}} \DecValTok{100} \OperatorTok{/}\NormalTok{ (}\DecValTok{1} \OperatorTok{+}\NormalTok{ gain }\OperatorTok{/}\NormalTok{ perte)}
\NormalTok{jours }\OperatorTok{=}\NormalTok{ pd.get\_dummies(d[}\StringTok{"Date"}\NormalTok{].dt.dayofweek, prefix}\OperatorTok{=}\StringTok{"jour"}\NormalTok{, dtype}\OperatorTok{=}\BuiltInTok{int}\NormalTok{)   }\CommentTok{\# encodage one{-}hot}
\NormalTok{X }\OperatorTok{=}\NormalTok{ pd.concat([X, jours], axis}\OperatorTok{=}\DecValTok{1}\NormalTok{)}

\CommentTok{\# Cible : le rendement de LA SÉANCE SUIVANTE est{-}il positif ?  (1 = hausse, 0 = baisse ou nul)}
\NormalTok{y }\OperatorTok{=}\NormalTok{ (r\_.shift(}\OperatorTok{{-}}\DecValTok{1}\NormalTok{) }\OperatorTok{\textgreater{}} \DecValTok{0}\NormalTok{).astype(}\BuiltInTok{float}\NormalTok{)}
\NormalTok{y[r\_.shift(}\OperatorTok{{-}}\DecValTok{1}\NormalTok{).isna()] }\OperatorTok{=}\NormalTok{ np.nan}

\NormalTok{base }\OperatorTok{=}\NormalTok{ pd.concat([d[[}\StringTok{"Date"}\NormalTok{]], X, y.rename(}\StringTok{"cible"}\NormalTok{)], axis}\OperatorTok{=}\DecValTok{1}\NormalTok{).dropna().reset\_index(drop}\OperatorTok{=}\VariableTok{True}\NormalTok{)}
\NormalTok{COL\_CONT }\OperatorTok{=}\NormalTok{ [c }\ControlFlowTok{for}\NormalTok{ c }\KeywordTok{in}\NormalTok{ X.columns }\ControlFlowTok{if} \KeywordTok{not}\NormalTok{ c.startswith(}\StringTok{"jour\_"}\NormalTok{)]}
\NormalTok{COL\_VAR }\OperatorTok{=} \BuiltInTok{list}\NormalTok{(X.columns)}
\BuiltInTok{print}\NormalTok{(}\StringTok{"Table finale :"}\NormalTok{, base.shape, }\StringTok{"| variables :"}\NormalTok{, }\BuiltInTok{len}\NormalTok{(COL\_VAR))}
\BuiltInTok{print}\NormalTok{(}\StringTok{"Période :"}\NormalTok{, base[}\StringTok{"Date"}\NormalTok{].}\BuiltInTok{min}\NormalTok{().date(), }\StringTok{"→"}\NormalTok{, base[}\StringTok{"Date"}\NormalTok{].}\BuiltInTok{max}\NormalTok{().date())}
\BuiltInTok{print}\NormalTok{(}\StringTok{"Équilibre des classes (part de hausses) :"}\NormalTok{, }\BuiltInTok{round}\NormalTok{(base[}\StringTok{"cible"}\NormalTok{].mean(), }\DecValTok{4}\NormalTok{))}
\BuiltInTok{print}\NormalTok{(}\StringTok{"Fenêtre perdue au démarrage (calcul des indicateurs) :"}\NormalTok{, }\BuiltInTok{len}\NormalTok{(d) }\OperatorTok{{-}} \BuiltInTok{len}\NormalTok{(base), }\StringTok{"lignes"}\NormalTok{)}
\end{Highlighting}
\end{Shaded}

\begin{Shaded}
\begin{Highlighting}[]
\NormalTok{Table finale : (1938, 21) | variables : 19}
\NormalTok{Période : 2018{-}03{-}13 → 2025{-}12{-}30}
\NormalTok{Équilibre des classes (part de hausses) : 0.5263}
\NormalTok{Fenêtre perdue au démarrage (calcul des indicateurs) : 50 lignes}
\end{Highlighting}
\end{Shaded}

\begin{Shaded}
\begin{Highlighting}[]
\NormalTok{n }\OperatorTok{=} \BuiltInTok{len}\NormalTok{(base)}
\NormalTok{i1, i2 }\OperatorTok{=} \BuiltInTok{int}\NormalTok{(n }\OperatorTok{*} \FloatTok{0.70}\NormalTok{), }\BuiltInTok{int}\NormalTok{(n }\OperatorTok{*} \FloatTok{0.85}\NormalTok{)}
\NormalTok{tr, va, te }\OperatorTok{=}\NormalTok{ base.iloc[:i1], base.iloc[i1:i2], base.iloc[i2:]}
\ControlFlowTok{for}\NormalTok{ nom, part }\KeywordTok{in}\NormalTok{ [(}\StringTok{"train"}\NormalTok{, tr), (}\StringTok{"validation"}\NormalTok{, va), (}\StringTok{"test"}\NormalTok{, te)]:}
    \BuiltInTok{print}\NormalTok{(}\SpecialStringTok{f"}\SpecialCharTok{\{}\NormalTok{nom}\SpecialCharTok{:10s\}}\SpecialStringTok{ }\SpecialCharTok{\{}\BuiltInTok{len}\NormalTok{(part)}\SpecialCharTok{:5d\}}\SpecialStringTok{ lignes  }\SpecialCharTok{\{}\NormalTok{part[}\StringTok{\textquotesingle{}Date\textquotesingle{}}\NormalTok{]}\SpecialCharTok{.}\BuiltInTok{min}\NormalTok{()}\SpecialCharTok{.}\NormalTok{date()}\SpecialCharTok{\}}\SpecialStringTok{ → }\SpecialCharTok{\{}\NormalTok{part[}\StringTok{\textquotesingle{}Date\textquotesingle{}}\NormalTok{]}\SpecialCharTok{.}\BuiltInTok{max}\NormalTok{()}\SpecialCharTok{.}\NormalTok{date()}\SpecialCharTok{\}}\SpecialStringTok{  part de hausses = }\SpecialCharTok{\{}\NormalTok{part[}\StringTok{\textquotesingle{}cible\textquotesingle{}}\NormalTok{]}\SpecialCharTok{.}\NormalTok{mean()}\SpecialCharTok{:.3f\}}\SpecialStringTok{"}\NormalTok{)}
\NormalTok{Xtr, ytr }\OperatorTok{=}\NormalTok{ tr[COL\_VAR], tr[}\StringTok{"cible"}\NormalTok{].astype(}\BuiltInTok{int}\NormalTok{)}
\NormalTok{Xva, yva }\OperatorTok{=}\NormalTok{ va[COL\_VAR], va[}\StringTok{"cible"}\NormalTok{].astype(}\BuiltInTok{int}\NormalTok{)}
\NormalTok{Xte, yte }\OperatorTok{=}\NormalTok{ te[COL\_VAR], te[}\StringTok{"cible"}\NormalTok{].astype(}\BuiltInTok{int}\NormalTok{)}
\CommentTok{\# La cible de la ligne t utilise le rendement t+1 ; la dernière ligne de chaque bloc partage donc}
\CommentTok{\# une information avec la première du bloc suivant : on neutralise cette fuite avec un embargo d\textquotesingle{}1 ligne.}
\NormalTok{Xva, yva }\OperatorTok{=}\NormalTok{ Xva.iloc[}\DecValTok{1}\NormalTok{:], yva.iloc[}\DecValTok{1}\NormalTok{:]}
\NormalTok{base.to\_csv(DOSSIER\_PROC }\OperatorTok{/} \StringTok{"base\_modelisation.csv"}\NormalTok{, index}\OperatorTok{=}\VariableTok{False}\NormalTok{)   }\CommentTok{\# table préparée (variables + cible)}
\NormalTok{Xte, yte }\OperatorTok{=}\NormalTok{ Xte.iloc[}\DecValTok{1}\NormalTok{:], yte.iloc[}\DecValTok{1}\NormalTok{:]}
\end{Highlighting}
\end{Shaded}

\begin{Shaded}
\begin{Highlighting}[]
\NormalTok{train       1356 lignes  2018{-}03{-}13 → 2023{-}08{-}22  part de hausses = 0.507}
\NormalTok{validation   291 lignes  2023{-}08{-}23 → 2024{-}10{-}25  part de hausses = 0.553}
\NormalTok{test         291 lignes  2024{-}10{-}28 → 2025{-}12{-}30  part de hausses = 0.588}
\end{Highlighting}
\end{Shaded}

\hypertarget{algorithmes}{%
\subsection{5. Algorithmes}\label{algorithmes}}

\begin{longtable}[]{@{}
  >{\raggedright\arraybackslash}p{(\columnwidth - 4\tabcolsep) * \real{0.3333}}
  >{\raggedright\arraybackslash}p{(\columnwidth - 4\tabcolsep) * \real{0.3333}}
  >{\raggedright\arraybackslash}p{(\columnwidth - 4\tabcolsep) * \real{0.3333}}@{}}
\toprule\noalign{}
\begin{minipage}[b]{\linewidth}\raggedright
Modèle
\end{minipage} & \begin{minipage}[b]{\linewidth}\raggedright
Rôle
\end{minipage} & \begin{minipage}[b]{\linewidth}\raggedright
Pourquoi
\end{minipage} \\
\midrule\noalign{}
\endhead
\bottomrule\noalign{}
\endlastfoot
\textbf{Référence « marche aléatoire »} & Adversaire principal & Sans
information exploitable, le meilleur pari est la classe majoritaire du
train (dérive positive). \\
Référence tirage pondéré & Contrôle & Prédire au hasard selon les
fréquences du train. \\
Référence persistance & Contrôle naïf & « Demain aura le même signe
qu'aujourd'hui » : teste directement l'autocorrélation observée. \\
\textbf{Régression logistique} (L2) & Modèle linéaire, interprétable &
Variables standardisées ; coefficient \texttt{C} choisi sur la
validation. \\
\textbf{Forêt aléatoire} & Modèle non linéaire robuste & Profondeur
limitée (2--5) et feuilles minimales (20/50) pour limiter le
surapprentissage sur un signal faible. \\
\textbf{Gradient boosting} (\texttt{HistGradientBoosting}) & Modèle non
linéaire séquentiel & Arbres peu profonds, faible taux
d'apprentissage. \\
\end{longtable}

\hypertarget{protocole}{%
\subsection{6. Protocole}\label{protocole}}

\begin{enumerate}
\def\labelenumi{\arabic{enumi}.}
\tightlist
\item
  \textbf{Hyperparamètres} choisis sur la \textbf{validation} (critère :
  précision équilibrée, adaptée à des classes légèrement
  déséquilibrées), jamais sur le test.
\item
  \textbf{Réentraînement} de chaque modèle sur train + validation, puis
  \textbf{évaluation unique} sur le test (graine 42 partout).
\item
  \textbf{Métriques} : précision, précision équilibrée, F1, MCC, AUC.
\item
  \textbf{Significativité} : test binomial exact (contre 50 \% et contre
  la référence majoritaire), intervalle de Wilson à 95 \%, test de
  McNemar exact (modèle contre référence).
\item
  \textbf{Robustesse} : validation glissante \emph{walk-forward}
  (2021-2025, réentraînement tous les 63 séances, fenêtre croissante),
  car 290 séances de test donnent un intervalle de ±5,6 points environ.
\item
  \textbf{Lecture économique} : stratégie « long si hausse prédite,
  sinon liquidités » avec un coût de 0,10 \% par changement de position
  (hypothèse), comparée à l'achat-conservation.
\end{enumerate}

\hypertarget{choix-des-hyperparamuxe8tres-validation}{%
\subsubsection{6.1 Choix des hyperparamètres
(validation)}\label{choix-des-hyperparamuxe8tres-validation}}

\begin{Shaded}
\begin{Highlighting}[]
\KeywordTok{def}\NormalTok{ preproc(echelle):}
    \CommentTok{\# La mise à l\textquotesingle{}échelle est apprise sur le TRAIN uniquement (dans le pipeline), jamais sur val/test.}
    \ControlFlowTok{return}\NormalTok{ ColumnTransformer([(}\StringTok{"sc"}\NormalTok{, StandardScaler() }\ControlFlowTok{if}\NormalTok{ echelle }\ControlFlowTok{else} \StringTok{"passthrough"}\NormalTok{, COL\_CONT)],}
\NormalTok{                             remainder}\OperatorTok{=}\StringTok{"passthrough"}\NormalTok{)}

\KeywordTok{def}\NormalTok{ metriques(y\_vrai, y\_pred, proba}\OperatorTok{=}\VariableTok{None}\NormalTok{):}
\NormalTok{    m }\OperatorTok{=}\NormalTok{ \{}\StringTok{"accuracy"}\NormalTok{: accuracy\_score(y\_vrai, y\_pred),}
         \StringTok{"acc\_equilibree"}\NormalTok{: balanced\_accuracy\_score(y\_vrai, y\_pred),}
         \StringTok{"F1"}\NormalTok{: f1\_score(y\_vrai, y\_pred, zero\_division}\OperatorTok{=}\DecValTok{0}\NormalTok{),}
         \StringTok{"MCC"}\NormalTok{: matthews\_corrcoef(y\_vrai, y\_pred)\}}
\NormalTok{    m[}\StringTok{"AUC"}\NormalTok{] }\OperatorTok{=}\NormalTok{ roc\_auc\_score(y\_vrai, proba) }\ControlFlowTok{if}\NormalTok{ proba }\KeywordTok{is} \KeywordTok{not} \VariableTok{None} \ControlFlowTok{else}\NormalTok{ np.nan}
    \ControlFlowTok{return}\NormalTok{ m}

\NormalTok{grilles }\OperatorTok{=}\NormalTok{ \{}
    \StringTok{"Régression logistique"}\NormalTok{: (}\KeywordTok{lambda} \OperatorTok{**}\NormalTok{p: Pipeline([(}\StringTok{"pre"}\NormalTok{, preproc(}\VariableTok{True}\NormalTok{)),}
\NormalTok{                              (}\StringTok{"m"}\NormalTok{, LogisticRegression(max\_iter}\OperatorTok{=}\DecValTok{2000}\NormalTok{, random\_state}\OperatorTok{=}\NormalTok{GRAINE, }\OperatorTok{**}\NormalTok{p))]),}
\NormalTok{                              [\{}\StringTok{"C"}\NormalTok{: c\} }\ControlFlowTok{for}\NormalTok{ c }\KeywordTok{in}\NormalTok{ [}\FloatTok{0.001}\NormalTok{, }\FloatTok{0.01}\NormalTok{, }\FloatTok{0.1}\NormalTok{, }\DecValTok{1}\NormalTok{, }\DecValTok{10}\NormalTok{]]),}
    \StringTok{"Forêt aléatoire"}\NormalTok{: (}\KeywordTok{lambda} \OperatorTok{**}\NormalTok{p: Pipeline([(}\StringTok{"pre"}\NormalTok{, preproc(}\VariableTok{False}\NormalTok{)),}
\NormalTok{                        (}\StringTok{"m"}\NormalTok{, RandomForestClassifier(n\_estimators}\OperatorTok{=}\DecValTok{300}\NormalTok{, random\_state}\OperatorTok{=}\NormalTok{GRAINE, n\_jobs}\OperatorTok{={-}}\DecValTok{1}\NormalTok{, }\OperatorTok{**}\NormalTok{p))]),}
\NormalTok{                        [\{}\StringTok{"max\_depth"}\NormalTok{: d\_, }\StringTok{"min\_samples\_leaf"}\NormalTok{: l\} }\ControlFlowTok{for}\NormalTok{ d\_ }\KeywordTok{in}\NormalTok{ [}\DecValTok{2}\NormalTok{, }\DecValTok{3}\NormalTok{, }\DecValTok{5}\NormalTok{] }\ControlFlowTok{for}\NormalTok{ l }\KeywordTok{in}\NormalTok{ [}\DecValTok{20}\NormalTok{, }\DecValTok{50}\NormalTok{]]),}
    \StringTok{"Gradient boosting"}\NormalTok{: (}\KeywordTok{lambda} \OperatorTok{**}\NormalTok{p: Pipeline([(}\StringTok{"pre"}\NormalTok{, preproc(}\VariableTok{False}\NormalTok{)),}
\NormalTok{                          (}\StringTok{"m"}\NormalTok{, HistGradientBoostingClassifier(random\_state}\OperatorTok{=}\NormalTok{GRAINE, }\OperatorTok{**}\NormalTok{p))]),}
\NormalTok{                          [\{}\StringTok{"max\_depth"}\NormalTok{: d\_, }\StringTok{"learning\_rate"}\NormalTok{: lr, }\StringTok{"max\_iter"}\NormalTok{: it, }\StringTok{"min\_samples\_leaf"}\NormalTok{: }\DecValTok{30}\NormalTok{\}}
                           \ControlFlowTok{for}\NormalTok{ d\_ }\KeywordTok{in}\NormalTok{ [}\DecValTok{2}\NormalTok{, }\DecValTok{3}\NormalTok{] }\ControlFlowTok{for}\NormalTok{ lr }\KeywordTok{in}\NormalTok{ [}\FloatTok{0.03}\NormalTok{, }\FloatTok{0.1}\NormalTok{] }\ControlFlowTok{for}\NormalTok{ it }\KeywordTok{in}\NormalTok{ [}\DecValTok{50}\NormalTok{, }\DecValTok{150}\NormalTok{]]),}
\NormalTok{\}}
\NormalTok{choix }\OperatorTok{=}\NormalTok{ \{\}}
\NormalTok{lignes }\OperatorTok{=}\NormalTok{ []}
\ControlFlowTok{for}\NormalTok{ nom, (fabrique, grille) }\KeywordTok{in}\NormalTok{ grilles.items():}
\NormalTok{    meilleur }\OperatorTok{=} \VariableTok{None}
    \ControlFlowTok{for}\NormalTok{ p }\KeywordTok{in}\NormalTok{ grille:}
\NormalTok{        mod }\OperatorTok{=}\NormalTok{ fabrique(}\OperatorTok{**}\NormalTok{p).fit(Xtr, ytr)}
\NormalTok{        proba }\OperatorTok{=}\NormalTok{ mod.predict\_proba(Xva)[:, }\DecValTok{1}\NormalTok{]}
\NormalTok{        pred }\OperatorTok{=}\NormalTok{ (proba }\OperatorTok{\textgreater{}=} \FloatTok{0.5}\NormalTok{).astype(}\BuiltInTok{int}\NormalTok{)}
\NormalTok{        score }\OperatorTok{=}\NormalTok{ balanced\_accuracy\_score(yva, pred)}
\NormalTok{        lignes.append(\{}\StringTok{"modèle"}\NormalTok{: nom, }\StringTok{"hyperparamètres"}\NormalTok{: }\BuiltInTok{str}\NormalTok{(p), }\StringTok{"acc\_équilibrée\_val"}\NormalTok{: }\BuiltInTok{round}\NormalTok{(score, }\DecValTok{4}\NormalTok{),}
                       \StringTok{"AUC\_val"}\NormalTok{: }\BuiltInTok{round}\NormalTok{(roc\_auc\_score(yva, proba), }\DecValTok{4}\NormalTok{)\})}
        \ControlFlowTok{if}\NormalTok{ meilleur }\KeywordTok{is} \VariableTok{None} \KeywordTok{or}\NormalTok{ score }\OperatorTok{\textgreater{}}\NormalTok{ meilleur[}\DecValTok{0}\NormalTok{]:}
\NormalTok{            meilleur }\OperatorTok{=}\NormalTok{ (score, p)}
\NormalTok{    choix[nom] }\OperatorTok{=}\NormalTok{ meilleur[}\DecValTok{1}\NormalTok{]}
\BuiltInTok{print}\NormalTok{(pd.DataFrame(lignes).sort\_values([}\StringTok{"modèle"}\NormalTok{, }\StringTok{"acc\_équilibrée\_val"}\NormalTok{], ascending}\OperatorTok{=}\NormalTok{[}\VariableTok{True}\NormalTok{, }\VariableTok{False}\NormalTok{]).to\_string(index}\OperatorTok{=}\VariableTok{False}\NormalTok{))}
\BuiltInTok{print}\NormalTok{(}\StringTok{"}\CharTok{\textbackslash{}n}\StringTok{Hyperparamètres retenus (validation) :"}\NormalTok{)}
\ControlFlowTok{for}\NormalTok{ k, v }\KeywordTok{in}\NormalTok{ choix.items(): }\BuiltInTok{print}\NormalTok{(}\StringTok{" {-}"}\NormalTok{, k, v)}
\end{Highlighting}
\end{Shaded}

\begin{Shaded}
\begin{Highlighting}[]
\NormalTok{               modèle                                                                  hyperparamètres  acc\_équilibrée\_val  AUC\_val}
\NormalTok{      Forêt aléatoire                                         \{\textquotesingle{}max\_depth\textquotesingle{}: 5, \textquotesingle{}min\_samples\_leaf\textquotesingle{}: 20\}              0.5517   0.5425}
\NormalTok{      Forêt aléatoire                                         \{\textquotesingle{}max\_depth\textquotesingle{}: 5, \textquotesingle{}min\_samples\_leaf\textquotesingle{}: 50\}              0.5416   0.5453}
\NormalTok{      Forêt aléatoire                                         \{\textquotesingle{}max\_depth\textquotesingle{}: 3, \textquotesingle{}min\_samples\_leaf\textquotesingle{}: 20\}              0.5389   0.5499}
\NormalTok{      Forêt aléatoire                                         \{\textquotesingle{}max\_depth\textquotesingle{}: 3, \textquotesingle{}min\_samples\_leaf\textquotesingle{}: 50\}              0.5334   0.5488}
\NormalTok{      Forêt aléatoire                                         \{\textquotesingle{}max\_depth\textquotesingle{}: 2, \textquotesingle{}min\_samples\_leaf\textquotesingle{}: 50\}              0.5315   0.5516}
\NormalTok{      Forêt aléatoire                                         \{\textquotesingle{}max\_depth\textquotesingle{}: 2, \textquotesingle{}min\_samples\_leaf\textquotesingle{}: 20\}              0.5130   0.5544}
\NormalTok{    Gradient boosting  \{\textquotesingle{}max\_depth\textquotesingle{}: 2, \textquotesingle{}learning\_rate\textquotesingle{}: 0.1, \textquotesingle{}max\_iter\textquotesingle{}: 150, \textquotesingle{}min\_samples\_leaf\textquotesingle{}: 30\}              0.5512   0.5662}
\NormalTok{    Gradient boosting  \{\textquotesingle{}max\_depth\textquotesingle{}: 2, \textquotesingle{}learning\_rate\textquotesingle{}: 0.03, \textquotesingle{}max\_iter\textquotesingle{}: 50, \textquotesingle{}min\_samples\_leaf\textquotesingle{}: 30\}              0.5368   0.5300}
\NormalTok{    Gradient boosting  \{\textquotesingle{}max\_depth\textquotesingle{}: 3, \textquotesingle{}learning\_rate\textquotesingle{}: 0.1, \textquotesingle{}max\_iter\textquotesingle{}: 150, \textquotesingle{}min\_samples\_leaf\textquotesingle{}: 30\}              0.5358   0.5258}
\NormalTok{    Gradient boosting  \{\textquotesingle{}max\_depth\textquotesingle{}: 3, \textquotesingle{}learning\_rate\textquotesingle{}: 0.03, \textquotesingle{}max\_iter\textquotesingle{}: 50, \textquotesingle{}min\_samples\_leaf\textquotesingle{}: 30\}              0.5161   0.5105}
\NormalTok{    Gradient boosting \{\textquotesingle{}max\_depth\textquotesingle{}: 3, \textquotesingle{}learning\_rate\textquotesingle{}: 0.03, \textquotesingle{}max\_iter\textquotesingle{}: 150, \textquotesingle{}min\_samples\_leaf\textquotesingle{}: 30\}              0.5139   0.5172}
\NormalTok{    Gradient boosting   \{\textquotesingle{}max\_depth\textquotesingle{}: 2, \textquotesingle{}learning\_rate\textquotesingle{}: 0.1, \textquotesingle{}max\_iter\textquotesingle{}: 50, \textquotesingle{}min\_samples\_leaf\textquotesingle{}: 30\}              0.5101   0.5250}
\NormalTok{    Gradient boosting \{\textquotesingle{}max\_depth\textquotesingle{}: 2, \textquotesingle{}learning\_rate\textquotesingle{}: 0.03, \textquotesingle{}max\_iter\textquotesingle{}: 150, \textquotesingle{}min\_samples\_leaf\textquotesingle{}: 30\}              0.5062   0.5308}
\NormalTok{    Gradient boosting   \{\textquotesingle{}max\_depth\textquotesingle{}: 3, \textquotesingle{}learning\_rate\textquotesingle{}: 0.1, \textquotesingle{}max\_iter\textquotesingle{}: 50, \textquotesingle{}min\_samples\_leaf\textquotesingle{}: 30\}              0.5053   0.5096}
\NormalTok{Régression logistique                                                                      \{\textquotesingle{}C\textquotesingle{}: 0.01\}              0.4964   0.5067}
\NormalTok{Régression logistique                                                                         \{\textquotesingle{}C\textquotesingle{}: 1\}              0.4887   0.4898}
\NormalTok{Régression logistique                                                                        \{\textquotesingle{}C\textquotesingle{}: 10\}              0.4887   0.4895}
\NormalTok{Régression logistique                                                                       \{\textquotesingle{}C\textquotesingle{}: 0.1\}              0.4873   0.4929}
\NormalTok{Régression logistique                                                                     \{\textquotesingle{}C\textquotesingle{}: 0.001\}              0.4839   0.5235}

\NormalTok{Hyperparamètres retenus (validation) :}
\NormalTok{ {-} Régression logistique \{\textquotesingle{}C\textquotesingle{}: 0.01\}}
\NormalTok{ {-} Forêt aléatoire \{\textquotesingle{}max\_depth\textquotesingle{}: 5, \textquotesingle{}min\_samples\_leaf\textquotesingle{}: 20\}}
\NormalTok{ {-} Gradient boosting \{\textquotesingle{}max\_depth\textquotesingle{}: 2, \textquotesingle{}learning\_rate\textquotesingle{}: 0.1, \textquotesingle{}max\_iter\textquotesingle{}: 150, \textquotesingle{}min\_samples\_leaf\textquotesingle{}: 30\}}
\end{Highlighting}
\end{Shaded}

\hypertarget{ruxe9sultats}{%
\subsection{7. Résultats}\label{ruxe9sultats}}

\hypertarget{test-chronologique-jamais-vu-pendant-la-suxe9lection}{%
\subsubsection{7.1 Test chronologique (jamais vu pendant la
sélection)}\label{test-chronologique-jamais-vu-pendant-la-suxe9lection}}

\begin{Shaded}
\begin{Highlighting}[]
\NormalTok{Xdev }\OperatorTok{=}\NormalTok{ pd.concat([Xtr, Xva])}\OperatorTok{;}\NormalTok{ ydev }\OperatorTok{=}\NormalTok{ pd.concat([ytr, yva])}
\NormalTok{resultats, predictions, probas }\OperatorTok{=}\NormalTok{ \{\}, \{\}, \{\}}

\CommentTok{\# {-}{-}{-} Références {-}{-}{-}}
\NormalTok{maj }\OperatorTok{=}\NormalTok{ DummyClassifier(strategy}\OperatorTok{=}\StringTok{"most\_frequent"}\NormalTok{).fit(Xdev, ydev)}
\NormalTok{alea }\OperatorTok{=}\NormalTok{ DummyClassifier(strategy}\OperatorTok{=}\StringTok{"stratified"}\NormalTok{, random\_state}\OperatorTok{=}\NormalTok{GRAINE).fit(Xdev, ydev)}
\NormalTok{predictions[}\StringTok{"Réf. marche aléatoire (classe majoritaire)"}\NormalTok{] }\OperatorTok{=}\NormalTok{ maj.predict(Xte)}
\NormalTok{predictions[}\StringTok{"Réf. tirage aléatoire pondéré"}\NormalTok{] }\OperatorTok{=}\NormalTok{ alea.predict(Xte)}
\CommentTok{\# Persistance : « demain = signe d\textquotesingle{}aujourd\textquotesingle{}hui » (rendement du jour t = r\_lag0)}
\NormalTok{predictions[}\StringTok{"Réf. persistance (signe de la veille)"}\NormalTok{] }\OperatorTok{=}\NormalTok{ (Xte[}\StringTok{"r\_lag0"}\NormalTok{] }\OperatorTok{\textgreater{}} \DecValTok{0}\NormalTok{).astype(}\BuiltInTok{int}\NormalTok{).values}
\NormalTok{probas[}\StringTok{"Réf. persistance (signe de la veille)"}\NormalTok{] }\OperatorTok{=}\NormalTok{ Xte[}\StringTok{"r\_lag0"}\NormalTok{].values}

\CommentTok{\# {-}{-}{-} Modèles (réentraînés sur train+validation avec les hyperparamètres choisis) {-}{-}{-}}
\ControlFlowTok{for}\NormalTok{ nom, (fabrique, \_) }\KeywordTok{in}\NormalTok{ grilles.items():}
\NormalTok{    mod }\OperatorTok{=}\NormalTok{ fabrique(}\OperatorTok{**}\NormalTok{choix[nom]).fit(Xdev, ydev)}
\NormalTok{    probas[nom] }\OperatorTok{=}\NormalTok{ mod.predict\_proba(Xte)[:, }\DecValTok{1}\NormalTok{]}
\NormalTok{    predictions[nom] }\OperatorTok{=}\NormalTok{ (probas[nom] }\OperatorTok{\textgreater{}=} \FloatTok{0.5}\NormalTok{).astype(}\BuiltInTok{int}\NormalTok{)}

\ControlFlowTok{for}\NormalTok{ nom, pred }\KeywordTok{in}\NormalTok{ predictions.items():}
\NormalTok{    resultats[nom] }\OperatorTok{=}\NormalTok{ metriques(yte, pred, probas.get(nom))}
\NormalTok{tab }\OperatorTok{=}\NormalTok{ pd.DataFrame(resultats).T.}\BuiltInTok{round}\NormalTok{(}\DecValTok{4}\NormalTok{)}
\NormalTok{display(tab)}
\BuiltInTok{print}\NormalTok{(}\StringTok{"}\CharTok{\textbackslash{}n}\StringTok{Test :"}\NormalTok{, }\BuiltInTok{len}\NormalTok{(yte), }\StringTok{"séances,"}\NormalTok{, Xte.index.size, }\StringTok{"| part de hausses :"}\NormalTok{, }\BuiltInTok{round}\NormalTok{(yte.mean(), }\DecValTok{4}\NormalTok{))}
\end{Highlighting}
\end{Shaded}

\begin{Shaded}
\begin{Highlighting}[]
\NormalTok{                                            accuracy  acc\_equilibree      F1     MCC     AUC}
\NormalTok{Réf. marche aléatoire (classe majoritaire)    0.5862          0.5000  0.7391  0.0000     NaN}
\NormalTok{Réf. tirage aléatoire pondéré                 0.4621          0.4603  0.5063 {-}0.0782     NaN}
\NormalTok{Réf. persistance (signe de la veille)         0.5793          0.5664  0.6412  0.1328  0.6018}
\NormalTok{Régression logistique                         0.5345          0.5159  0.6110  0.0322  0.5338}
\NormalTok{Forêt aléatoire                               0.5310          0.5179  0.5976  0.0357  0.5345}
\NormalTok{Gradient boosting                             0.5034          0.5103  0.5263  0.0203  0.5194}

\NormalTok{Test : 290 séances, 290 | part de hausses : 0.5862}
\end{Highlighting}
\end{Shaded}

\textbf{Lecture.} Sur le test, la référence « toujours hausse » atteint
\textbf{58,6 \%} car la période (oct. 2024 → déc. 2025) est fortement
haussière (58,8 \% de séances positives). Les trois modèles font moins
bien : régression logistique 53,4 \%, forêt 53,1 \%, gradient boosting
50,3 \%. Leur précision équilibrée (≈ 0,51-0,52), leur MCC (≈ 0,02-0,04)
et leur AUC (≈ 0,52-0,53) sont proches du hasard. La persistance obtient
la meilleure AUC (0,60) et le meilleur MCC (0,13).

\hypertarget{significativituxe9-statistique}{%
\subsubsection{7.2 Significativité
statistique}\label{significativituxe9-statistique}}

\begin{Shaded}
\begin{Highlighting}[]
\NormalTok{ref\_nom }\OperatorTok{=} \StringTok{"Réf. marche aléatoire (classe majoritaire)"}
\NormalTok{ref\_pred }\OperatorTok{=}\NormalTok{ predictions[ref\_nom]}
\NormalTok{p0 }\OperatorTok{=}\NormalTok{ accuracy\_score(yte, ref\_pred)}
\BuiltInTok{print}\NormalTok{(}\SpecialStringTok{f"Précision de la référence (toujours « hausse ») sur le test : }\SpecialCharTok{\{}\NormalTok{p0}\SpecialCharTok{:.4f\}}\CharTok{\textbackslash{}n}\SpecialStringTok{"}\NormalTok{)}
\NormalTok{lig }\OperatorTok{=}\NormalTok{ []}
\ControlFlowTok{for}\NormalTok{ nom }\KeywordTok{in}\NormalTok{ [}\StringTok{"Régression logistique"}\NormalTok{, }\StringTok{"Forêt aléatoire"}\NormalTok{, }\StringTok{"Gradient boosting"}\NormalTok{, }\StringTok{"Réf. persistance (signe de la veille)"}\NormalTok{]:}
\NormalTok{    pred }\OperatorTok{=}\NormalTok{ predictions[nom]}
\NormalTok{    k }\OperatorTok{=} \BuiltInTok{int}\NormalTok{((pred }\OperatorTok{==}\NormalTok{ yte.values).}\BuiltInTok{sum}\NormalTok{())}\OperatorTok{;}\NormalTok{ N }\OperatorTok{=} \BuiltInTok{len}\NormalTok{(yte)}
\NormalTok{    ic }\OperatorTok{=}\NormalTok{ stats.binomtest(k, N).proportion\_ci(}\FloatTok{0.95}\NormalTok{, method}\OperatorTok{=}\StringTok{"wilson"}\NormalTok{)}
\NormalTok{    p\_bin }\OperatorTok{=}\NormalTok{ stats.binomtest(k, N, }\FloatTok{0.5}\NormalTok{, alternative}\OperatorTok{=}\StringTok{"greater"}\NormalTok{).pvalue          }\CommentTok{\# vs hasard pur 50 \%}
\NormalTok{    p\_ref }\OperatorTok{=}\NormalTok{ stats.binomtest(k, N, p0, alternative}\OperatorTok{=}\StringTok{"greater"}\NormalTok{).pvalue           }\CommentTok{\# vs référence majoritaire}
\NormalTok{    b }\OperatorTok{=} \BuiltInTok{int}\NormalTok{(((pred }\OperatorTok{==}\NormalTok{ yte.values) }\OperatorTok{\&}\NormalTok{ (ref\_pred }\OperatorTok{!=}\NormalTok{ yte.values)).}\BuiltInTok{sum}\NormalTok{())          }\CommentTok{\# modèle juste, réf. fausse}
\NormalTok{    c }\OperatorTok{=} \BuiltInTok{int}\NormalTok{(((pred }\OperatorTok{!=}\NormalTok{ yte.values) }\OperatorTok{\&}\NormalTok{ (ref\_pred }\OperatorTok{==}\NormalTok{ yte.values)).}\BuiltInTok{sum}\NormalTok{())          }\CommentTok{\# modèle faux, réf. juste}
\NormalTok{    p\_mc }\OperatorTok{=}\NormalTok{ stats.binomtest(b, b }\OperatorTok{+}\NormalTok{ c, }\FloatTok{0.5}\NormalTok{).pvalue }\ControlFlowTok{if}\NormalTok{ b }\OperatorTok{+}\NormalTok{ c }\OperatorTok{\textgreater{}} \DecValTok{0} \ControlFlowTok{else}\NormalTok{ np.nan     }\CommentTok{\# McNemar exact (bilatéral)}
\NormalTok{    lig.append(\{}\StringTok{"modèle"}\NormalTok{: nom, }\StringTok{"précision"}\NormalTok{: }\BuiltInTok{round}\NormalTok{(k }\OperatorTok{/}\NormalTok{ N, }\DecValTok{4}\NormalTok{), }\StringTok{"IC95 bas"}\NormalTok{: }\BuiltInTok{round}\NormalTok{(ic.low, }\DecValTok{3}\NormalTok{), }\StringTok{"IC95 haut"}\NormalTok{: }\BuiltInTok{round}\NormalTok{(ic.high, }\DecValTok{3}\NormalTok{),}
                \StringTok{"p (\textgreater{}50 \%)"}\NormalTok{: }\BuiltInTok{round}\NormalTok{(p\_bin, }\DecValTok{3}\NormalTok{), }\StringTok{"p (\textgreater{} réf. majoritaire)"}\NormalTok{: }\BuiltInTok{round}\NormalTok{(p\_ref, }\DecValTok{3}\NormalTok{), }\StringTok{"p McNemar vs réf."}\NormalTok{: }\BuiltInTok{round}\NormalTok{(p\_mc, }\DecValTok{3}\NormalTok{)\})}
\BuiltInTok{print}\NormalTok{(pd.DataFrame(lig).to\_string(index}\OperatorTok{=}\VariableTok{False}\NormalTok{))}
\end{Highlighting}
\end{Shaded}

\begin{Shaded}
\begin{Highlighting}[]
\NormalTok{Précision de la référence (toujours « hausse ») sur le test : 0.5862}

\NormalTok{                               modèle  précision  IC95 bas  IC95 haut  p (\textgreater{}50 \%)  p (\textgreater{} réf. majoritaire)  p McNemar vs réf.}
\NormalTok{                Régression logistique     0.5345     0.477      0.591      0.132                   0.967              0.188}
\NormalTok{                      Forêt aléatoire     0.5310     0.474      0.588      0.159                   0.975              0.174}
\NormalTok{                    Gradient boosting     0.5034     0.446      0.561      0.477                   0.998              0.065}
\NormalTok{Réf. persistance (signe de la veille)     0.5793     0.522      0.635      0.004                   0.618              0.927}
\end{Highlighting}
\end{Shaded}

\textbf{Lecture.} Aucun modèle ne dépasse la référence majoritaire (p ≈
0,97 à 1,00). Aucun n'est non plus significativement meilleur que 50 \%
(p = 0,13 ; 0,16 ; 0,48). Seule la persistance bat 50 \% (p = 0,004)
mais pas la référence majoritaire (p = 0,62). Les écarts avec la
référence ne sont pas significatifs au sens de McNemar (p ≥ 0,065).

\hypertarget{robustesse-validation-glissante-walk-forward-2021-2025}{%
\subsubsection{7.3 Robustesse : validation glissante (walk-forward,
2021-2025)}\label{robustesse-validation-glissante-walk-forward-2021-2025}}

\begin{Shaded}
\begin{Highlighting}[]
\CommentTok{\# Robustesse : validation « glissante » (walk{-}forward) sur 2021{-}2025, réentraînement tous les 63 séances,}
\CommentTok{\# fenêtre croissante, hyperparamètres figés (choisis ci{-}dessus).}
\NormalTok{DEB }\OperatorTok{=}\NormalTok{ base.index[base[}\StringTok{"Date"}\NormalTok{] }\OperatorTok{\textgreater{}=} \StringTok{"2021{-}01{-}01"}\NormalTok{][}\DecValTok{0}\NormalTok{]}
\NormalTok{PAS }\OperatorTok{=} \DecValTok{63}
\NormalTok{res\_wf }\OperatorTok{=}\NormalTok{ \{nom: [] }\ControlFlowTok{for}\NormalTok{ nom }\KeywordTok{in} \BuiltInTok{list}\NormalTok{(grilles) }\OperatorTok{+}\NormalTok{ [}\StringTok{"Réf. majoritaire"}\NormalTok{, }\StringTok{"Persistance"}\NormalTok{]\}}
\NormalTok{vrai\_wf }\OperatorTok{=}\NormalTok{ []}
\ControlFlowTok{for}\NormalTok{ debut }\KeywordTok{in} \BuiltInTok{range}\NormalTok{(DEB, }\BuiltInTok{len}\NormalTok{(base) }\OperatorTok{{-}} \DecValTok{1}\NormalTok{, PAS):}
\NormalTok{    fin }\OperatorTok{=} \BuiltInTok{min}\NormalTok{(debut }\OperatorTok{+}\NormalTok{ PAS, }\BuiltInTok{len}\NormalTok{(base))}
\NormalTok{    app }\OperatorTok{=}\NormalTok{ base.iloc[:debut }\OperatorTok{{-}} \DecValTok{1}\NormalTok{]                       }\CommentTok{\# embargo d\textquotesingle{}1 ligne}
\NormalTok{    tst }\OperatorTok{=}\NormalTok{ base.iloc[debut:fin]}
\NormalTok{    vrai\_wf.append(tst[}\StringTok{"cible"}\NormalTok{].astype(}\BuiltInTok{int}\NormalTok{).values)}
    \ControlFlowTok{for}\NormalTok{ nom, (fabrique, \_) }\KeywordTok{in}\NormalTok{ grilles.items():}
\NormalTok{        res\_wf[nom].append(fabrique(}\OperatorTok{**}\NormalTok{choix[nom]).fit(app[COL\_VAR], app[}\StringTok{"cible"}\NormalTok{].astype(}\BuiltInTok{int}\NormalTok{)).predict(tst[COL\_VAR]))}
\NormalTok{    res\_wf[}\StringTok{"Réf. majoritaire"}\NormalTok{].append(np.full(}\BuiltInTok{len}\NormalTok{(tst), }\BuiltInTok{int}\NormalTok{(app[}\StringTok{"cible"}\NormalTok{].mean() }\OperatorTok{\textgreater{}=} \FloatTok{0.5}\NormalTok{)))}
\NormalTok{    res\_wf[}\StringTok{"Persistance"}\NormalTok{].append((tst[}\StringTok{"r\_lag0"}\NormalTok{] }\OperatorTok{\textgreater{}} \DecValTok{0}\NormalTok{).astype(}\BuiltInTok{int}\NormalTok{).values)}
\NormalTok{vrai\_wf }\OperatorTok{=}\NormalTok{ np.concatenate(vrai\_wf)}
\NormalTok{lig }\OperatorTok{=}\NormalTok{ []}
\NormalTok{base\_maj }\OperatorTok{=}\NormalTok{ np.concatenate(res\_wf[}\StringTok{"Réf. majoritaire"}\NormalTok{])}
\ControlFlowTok{for}\NormalTok{ nom, morceaux }\KeywordTok{in}\NormalTok{ res\_wf.items():}
\NormalTok{    pred }\OperatorTok{=}\NormalTok{ np.concatenate(morceaux)}
\NormalTok{    k }\OperatorTok{=} \BuiltInTok{int}\NormalTok{((pred }\OperatorTok{==}\NormalTok{ vrai\_wf).}\BuiltInTok{sum}\NormalTok{())}\OperatorTok{;}\NormalTok{ N }\OperatorTok{=} \BuiltInTok{len}\NormalTok{(vrai\_wf)}
\NormalTok{    b }\OperatorTok{=} \BuiltInTok{int}\NormalTok{(((pred }\OperatorTok{==}\NormalTok{ vrai\_wf) }\OperatorTok{\&}\NormalTok{ (base\_maj }\OperatorTok{!=}\NormalTok{ vrai\_wf)).}\BuiltInTok{sum}\NormalTok{())}\OperatorTok{;}\NormalTok{ c }\OperatorTok{=} \BuiltInTok{int}\NormalTok{(((pred }\OperatorTok{!=}\NormalTok{ vrai\_wf) }\OperatorTok{\&}\NormalTok{ (base\_maj }\OperatorTok{==}\NormalTok{ vrai\_wf)).}\BuiltInTok{sum}\NormalTok{())}
\NormalTok{    lig.append(\{}\StringTok{"méthode"}\NormalTok{: nom, }\StringTok{"N"}\NormalTok{: N, }\StringTok{"précision"}\NormalTok{: }\BuiltInTok{round}\NormalTok{(k }\OperatorTok{/}\NormalTok{ N, }\DecValTok{4}\NormalTok{),}
                \StringTok{"acc\_équilibrée"}\NormalTok{: }\BuiltInTok{round}\NormalTok{(balanced\_accuracy\_score(vrai\_wf, pred), }\DecValTok{4}\NormalTok{),}
                \StringTok{"p (\textgreater{} 50 \%)"}\NormalTok{: }\BuiltInTok{round}\NormalTok{(stats.binomtest(k, N, }\FloatTok{0.5}\NormalTok{, alternative}\OperatorTok{=}\StringTok{"greater"}\NormalTok{).pvalue, }\DecValTok{4}\NormalTok{),}
                \StringTok{"p McNemar vs réf."}\NormalTok{: }\BuiltInTok{round}\NormalTok{(stats.binomtest(b, b }\OperatorTok{+}\NormalTok{ c, }\FloatTok{0.5}\NormalTok{).pvalue, }\DecValTok{4}\NormalTok{) }\ControlFlowTok{if}\NormalTok{ b }\OperatorTok{+}\NormalTok{ c }\ControlFlowTok{else}\NormalTok{ np.nan\})}
\BuiltInTok{print}\NormalTok{(pd.DataFrame(lig).to\_string(index}\OperatorTok{=}\VariableTok{False}\NormalTok{))}
\end{Highlighting}
\end{Shaded}

\begin{Shaded}
\begin{Highlighting}[]
\NormalTok{              méthode    N  précision  acc\_équilibrée  p (\textgreater{} 50 \%)  p McNemar vs réf.}
\NormalTok{Régression logistique 1244     0.5442          0.5335      0.0010             0.5626}
\NormalTok{      Forêt aléatoire 1244     0.5611          0.5547      0.0000             0.1437}
\NormalTok{    Gradient boosting 1244     0.5378          0.5397      0.0042             0.8441}
\NormalTok{     Réf. majoritaire 1244     0.5330          0.5005      0.0108                NaN}
\NormalTok{          Persistance 1244     0.5514          0.5475      0.0002             0.3614}
\end{Highlighting}
\end{Shaded}

\textbf{Lecture.} Sur 1 244 séances, les modèles font légèrement mieux
que la référence (53,3 \%) : forêt 56,1 \%, persistance 55,1 \%,
régression logistique 54,4 \%, gradient boosting 53,8 \%. Ils battent
clairement 50 \% mais cela reflète surtout la dérive haussière ;
\textbf{contre la référence majoritaire, aucune différence n'est
significative} (McNemar p = 0,14 pour la forêt, la meilleure).
Précaution : les hyperparamètres ont été choisis sur une période
(validation) qui fait partie de ces plis, ce qui avantage légèrement les
modèles.

\hypertarget{graphiques-de-ruxe9sultats}{%
\subsubsection{7.4 Graphiques de
résultats}\label{graphiques-de-ruxe9sultats}}

\begin{Shaded}
\begin{Highlighting}[]
\NormalTok{fig, ax }\OperatorTok{=}\NormalTok{ plt.subplots(}\DecValTok{1}\NormalTok{, }\DecValTok{2}\NormalTok{, figsize}\OperatorTok{=}\NormalTok{(}\DecValTok{13}\NormalTok{, }\FloatTok{4.5}\NormalTok{))}
\NormalTok{noms }\OperatorTok{=} \BuiltInTok{list}\NormalTok{(tab.index)}
\NormalTok{ax[}\DecValTok{0}\NormalTok{].barh(noms, tab[}\StringTok{"accuracy"}\NormalTok{], color}\OperatorTok{=}\NormalTok{[}\StringTok{"\#999"}\NormalTok{] }\OperatorTok{*} \DecValTok{3} \OperatorTok{+}\NormalTok{ [}\StringTok{"\#1f77b4"}\NormalTok{] }\OperatorTok{*} \DecValTok{3}\NormalTok{)}
\NormalTok{ax[}\DecValTok{0}\NormalTok{].axvline(}\FloatTok{0.5}\NormalTok{, color}\OperatorTok{=}\StringTok{"k"}\NormalTok{, ls}\OperatorTok{=}\StringTok{":"}\NormalTok{, label}\OperatorTok{=}\StringTok{"hasard pur (50 \%)"}\NormalTok{)}
\NormalTok{ax[}\DecValTok{0}\NormalTok{].axvline(p0, color}\OperatorTok{=}\StringTok{"r"}\NormalTok{, ls}\OperatorTok{=}\StringTok{"{-}{-}"}\NormalTok{, label}\OperatorTok{=}\StringTok{"référence majoritaire"}\NormalTok{)}
\NormalTok{ax[}\DecValTok{0}\NormalTok{].}\BuiltInTok{set}\NormalTok{(title}\OperatorTok{=}\StringTok{"Précision sur le jeu de test"}\NormalTok{, xlabel}\OperatorTok{=}\StringTok{"Précision (part de prévisions correctes)"}\NormalTok{, xlim}\OperatorTok{=}\NormalTok{(}\FloatTok{0.4}\NormalTok{, }\BuiltInTok{max}\NormalTok{(}\FloatTok{0.7}\NormalTok{, tab[}\StringTok{\textquotesingle{}accuracy\textquotesingle{}}\NormalTok{].}\BuiltInTok{max}\NormalTok{() }\OperatorTok{+} \FloatTok{.03}\NormalTok{)))}
\NormalTok{ax[}\DecValTok{0}\NormalTok{].legend(loc}\OperatorTok{=}\StringTok{"lower right"}\NormalTok{)}
\NormalTok{meilleur\_nom }\OperatorTok{=}\NormalTok{ tab.loc[[}\StringTok{"Régression logistique"}\NormalTok{, }\StringTok{"Forêt aléatoire"}\NormalTok{, }\StringTok{"Gradient boosting"}\NormalTok{], }\StringTok{"acc\_equilibree"}\NormalTok{].idxmax()}
\NormalTok{cm }\OperatorTok{=}\NormalTok{ confusion\_matrix(yte, predictions[meilleur\_nom])}
\NormalTok{ax[}\DecValTok{1}\NormalTok{].imshow(cm, cmap}\OperatorTok{=}\StringTok{"Blues"}\NormalTok{)}
\ControlFlowTok{for}\NormalTok{ (i, j), v }\KeywordTok{in}\NormalTok{ np.ndenumerate(cm): ax[}\DecValTok{1}\NormalTok{].text(j, i, v, ha}\OperatorTok{=}\StringTok{"center"}\NormalTok{, va}\OperatorTok{=}\StringTok{"center"}\NormalTok{, fontsize}\OperatorTok{=}\DecValTok{14}\NormalTok{)}
\NormalTok{ax[}\DecValTok{1}\NormalTok{].}\BuiltInTok{set}\NormalTok{(title}\OperatorTok{=}\SpecialStringTok{f"Matrice de confusion – }\SpecialCharTok{\{}\NormalTok{meilleur\_nom}\SpecialCharTok{\}}\SpecialStringTok{"}\NormalTok{, xlabel}\OperatorTok{=}\StringTok{"Classe prédite (0 = baisse, 1 = hausse)"}\NormalTok{, ylabel}\OperatorTok{=}\StringTok{"Classe réelle"}\NormalTok{,}
\NormalTok{          xticks}\OperatorTok{=}\NormalTok{[}\DecValTok{0}\NormalTok{, }\DecValTok{1}\NormalTok{], yticks}\OperatorTok{=}\NormalTok{[}\DecValTok{0}\NormalTok{, }\DecValTok{1}\NormalTok{])}\OperatorTok{;}\NormalTok{ ax[}\DecValTok{1}\NormalTok{].grid(}\VariableTok{False}\NormalTok{)}
\NormalTok{plt.tight\_layout()}\OperatorTok{;}\NormalTok{ plt.savefig(DOSSIER\_FIG }\OperatorTok{/} \StringTok{"02\_resultats\_test.png"}\NormalTok{, dpi}\OperatorTok{=}\DecValTok{150}\NormalTok{, bbox\_inches}\OperatorTok{=}\StringTok{"tight"}\NormalTok{)}\OperatorTok{;}\NormalTok{ plt.show()}
\end{Highlighting}
\end{Shaded}

\begin{figure}
\centering
\includegraphics[width=1\textwidth,height=\textheight]{../../figures/02_resultats_test.png}
\caption{Résultats sur le jeu de test : précision et matrice de
confusion}
\end{figure}

\hypertarget{variables-les-plus-influentes}{%
\subsubsection{7.5 Variables les plus
influentes}\label{variables-les-plus-influentes}}

\begin{Shaded}
\begin{Highlighting}[]
\NormalTok{rf }\OperatorTok{=}\NormalTok{ grilles[}\StringTok{"Forêt aléatoire"}\NormalTok{][}\DecValTok{0}\NormalTok{](}\OperatorTok{**}\NormalTok{choix[}\StringTok{"Forêt aléatoire"}\NormalTok{]).fit(Xdev, ydev)}
\NormalTok{noms\_var }\OperatorTok{=}\NormalTok{ COL\_CONT }\OperatorTok{+}\NormalTok{ [c }\ControlFlowTok{for}\NormalTok{ c }\KeywordTok{in}\NormalTok{ COL\_VAR }\ControlFlowTok{if}\NormalTok{ c }\KeywordTok{not} \KeywordTok{in}\NormalTok{ COL\_CONT]}
\NormalTok{imp }\OperatorTok{=}\NormalTok{ pd.Series(rf.named\_steps[}\StringTok{"m"}\NormalTok{].feature\_importances\_, index}\OperatorTok{=}\NormalTok{noms\_var).sort\_values()}
\NormalTok{lr }\OperatorTok{=}\NormalTok{ grilles[}\StringTok{"Régression logistique"}\NormalTok{][}\DecValTok{0}\NormalTok{](}\OperatorTok{**}\NormalTok{choix[}\StringTok{"Régression logistique"}\NormalTok{]).fit(Xdev, ydev)}
\NormalTok{coef }\OperatorTok{=}\NormalTok{ pd.Series(lr.named\_steps[}\StringTok{"m"}\NormalTok{].coef\_[}\DecValTok{0}\NormalTok{], index}\OperatorTok{=}\NormalTok{noms\_var).sort\_values()}
\NormalTok{fig, ax }\OperatorTok{=}\NormalTok{ plt.subplots(}\DecValTok{1}\NormalTok{, }\DecValTok{2}\NormalTok{, figsize}\OperatorTok{=}\NormalTok{(}\DecValTok{13}\NormalTok{, }\DecValTok{5}\NormalTok{))}
\NormalTok{imp.tail(}\DecValTok{10}\NormalTok{).plot.barh(ax}\OperatorTok{=}\NormalTok{ax[}\DecValTok{0}\NormalTok{])}\OperatorTok{;}\NormalTok{ ax[}\DecValTok{0}\NormalTok{].}\BuiltInTok{set}\NormalTok{(title}\OperatorTok{=}\StringTok{"Forêt aléatoire – 10 variables les plus importantes"}\NormalTok{, xlabel}\OperatorTok{=}\StringTok{"Importance (réduction d\textquotesingle{}impureté)"}\NormalTok{)}
\NormalTok{coef.reindex(coef.}\BuiltInTok{abs}\NormalTok{().sort\_values().tail(}\DecValTok{10}\NormalTok{).index).plot.barh(ax}\OperatorTok{=}\NormalTok{ax[}\DecValTok{1}\NormalTok{], color}\OperatorTok{=}\StringTok{"C1"}\NormalTok{)}
\NormalTok{ax[}\DecValTok{1}\NormalTok{].}\BuiltInTok{set}\NormalTok{(title}\OperatorTok{=}\StringTok{"Régression logistique – coefficients (variables standardisées)"}\NormalTok{, xlabel}\OperatorTok{=}\StringTok{"Coefficient"}\NormalTok{)}
\NormalTok{plt.tight\_layout()}\OperatorTok{;}\NormalTok{ plt.savefig(DOSSIER\_FIG }\OperatorTok{/} \StringTok{"03\_importance\_variables.png"}\NormalTok{, dpi}\OperatorTok{=}\DecValTok{150}\NormalTok{, bbox\_inches}\OperatorTok{=}\StringTok{"tight"}\NormalTok{)}\OperatorTok{;}\NormalTok{ plt.show()}
\end{Highlighting}
\end{Shaded}

\begin{figure}
\centering
\includegraphics[width=1\textwidth,height=\textheight]{../../figures/03_importance_variables.png}
\caption{Variables les plus influentes}
\end{figure}

\textbf{Lecture.} Dans les deux modèles, les variables dominantes sont
le \textbf{rendement du jour (\texttt{r\_lag0})} et le \textbf{RSI 14
jours} (coefficients positifs en régression logistique : continuation du
mouvement récent), ce qui rejoint l'autocorrélation positive de la
section 3. La forêt utilise aussi \texttt{moy\_20}, \texttt{r\_lag4} et
\texttt{moy\_5}. Ces importances décrivent le modèle, pas une causalité
ni un signal exploitable.

\hypertarget{lecture-uxe9conomique-backtest-simplifiuxe9-sur-le-test}{%
\subsubsection{7.6 Lecture économique (backtest simplifié sur le
test)}\label{lecture-uxe9conomique-backtest-simplifiuxe9-sur-le-test}}

\begin{Shaded}
\begin{Highlighting}[]
\NormalTok{rend\_te }\OperatorTok{=}\NormalTok{ te[}\StringTok{"Date"}\NormalTok{].iloc[}\DecValTok{1}\NormalTok{:].to\_frame()}
\NormalTok{rend\_te[}\StringTok{"r\_suivant"}\NormalTok{] }\OperatorTok{=}\NormalTok{ np.exp(np.log(df.set\_index(}\StringTok{"Date"}\NormalTok{)[}\StringTok{"Price"}\NormalTok{]).diff().shift(}\OperatorTok{{-}}\DecValTok{1}\NormalTok{).reindex(rend\_te[}\StringTok{"Date"}\NormalTok{]).values) }\OperatorTok{{-}} \DecValTok{1}
\NormalTok{cout }\OperatorTok{=} \FloatTok{0.001}                                           \CommentTok{\# 0,10 \% par changement de position (hypothèse)}
\NormalTok{lig }\OperatorTok{=}\NormalTok{ []}
\NormalTok{courbes }\OperatorTok{=}\NormalTok{ \{\}}
\ControlFlowTok{for}\NormalTok{ nom }\KeywordTok{in}\NormalTok{ [}\StringTok{"Régression logistique"}\NormalTok{, }\StringTok{"Forêt aléatoire"}\NormalTok{, }\StringTok{"Gradient boosting"}\NormalTok{]:}
\NormalTok{    pos }\OperatorTok{=}\NormalTok{ pd.Series(predictions[nom], index}\OperatorTok{=}\NormalTok{rend\_te.index)}
\NormalTok{    nb\_chg }\OperatorTok{=}\NormalTok{ pos.diff().}\BuiltInTok{abs}\NormalTok{().fillna(}\DecValTok{0}\NormalTok{)}
\NormalTok{    ret }\OperatorTok{=}\NormalTok{ pos }\OperatorTok{*}\NormalTok{ rend\_te[}\StringTok{"r\_suivant"}\NormalTok{] }\OperatorTok{{-}}\NormalTok{ nb\_chg }\OperatorTok{*}\NormalTok{ cout}
\NormalTok{    courbes[nom] }\OperatorTok{=}\NormalTok{ (}\DecValTok{1} \OperatorTok{+}\NormalTok{ ret).cumprod()}
\NormalTok{    lig.append(\{}\StringTok{"stratégie"}\NormalTok{: nom, }\StringTok{"rendement cumulé (\%)"}\NormalTok{: }\BuiltInTok{round}\NormalTok{((courbes[nom].iloc[}\OperatorTok{{-}}\DecValTok{1}\NormalTok{] }\OperatorTok{{-}} \DecValTok{1}\NormalTok{) }\OperatorTok{*} \DecValTok{100}\NormalTok{, }\DecValTok{2}\NormalTok{),}
                \StringTok{"Sharpe annualisé"}\NormalTok{: }\BuiltInTok{round}\NormalTok{(ret.mean() }\OperatorTok{/}\NormalTok{ ret.std() }\OperatorTok{*}\NormalTok{ np.sqrt(}\DecValTok{252}\NormalTok{), }\DecValTok{2}\NormalTok{) }\ControlFlowTok{if}\NormalTok{ ret.std() }\OperatorTok{\textgreater{}} \DecValTok{0} \ControlFlowTok{else}\NormalTok{ np.nan,}
                \StringTok{"nb de changements"}\NormalTok{: }\BuiltInTok{int}\NormalTok{(nb\_chg.}\BuiltInTok{sum}\NormalTok{())\})}
\NormalTok{bh }\OperatorTok{=}\NormalTok{ (}\DecValTok{1} \OperatorTok{+}\NormalTok{ rend\_te[}\StringTok{"r\_suivant"}\NormalTok{]).cumprod()}\OperatorTok{;}\NormalTok{ courbes[}\StringTok{"Acheter{-}conserver"}\NormalTok{] }\OperatorTok{=}\NormalTok{ bh}
\NormalTok{ret\_bh }\OperatorTok{=}\NormalTok{ rend\_te[}\StringTok{"r\_suivant"}\NormalTok{]}
\NormalTok{lig.append(\{}\StringTok{"stratégie"}\NormalTok{: }\StringTok{"Acheter{-}conserver"}\NormalTok{, }\StringTok{"rendement cumulé (\%)"}\NormalTok{: }\BuiltInTok{round}\NormalTok{((bh.iloc[}\OperatorTok{{-}}\DecValTok{1}\NormalTok{] }\OperatorTok{{-}} \DecValTok{1}\NormalTok{) }\OperatorTok{*} \DecValTok{100}\NormalTok{, }\DecValTok{2}\NormalTok{),}
            \StringTok{"Sharpe annualisé"}\NormalTok{: }\BuiltInTok{round}\NormalTok{(ret\_bh.mean() }\OperatorTok{/}\NormalTok{ ret\_bh.std() }\OperatorTok{*}\NormalTok{ np.sqrt(}\DecValTok{252}\NormalTok{), }\DecValTok{2}\NormalTok{), }\StringTok{"nb de changements"}\NormalTok{: }\DecValTok{0}\NormalTok{\})}
\BuiltInTok{print}\NormalTok{(pd.DataFrame(lig).to\_string(index}\OperatorTok{=}\VariableTok{False}\NormalTok{))}
\NormalTok{plt.figure(figsize}\OperatorTok{=}\NormalTok{(}\DecValTok{10}\NormalTok{, }\DecValTok{4}\NormalTok{))}
\ControlFlowTok{for}\NormalTok{ nom, c }\KeywordTok{in}\NormalTok{ courbes.items(): plt.plot(rend\_te[}\StringTok{"Date"}\NormalTok{].values, c.values, label}\OperatorTok{=}\NormalTok{nom)}
\NormalTok{plt.title(}\StringTok{"Valeur d\textquotesingle{}un portefeuille de 1 (test, coût 0,10 \% par changement)"}\NormalTok{)}\OperatorTok{;}\NormalTok{ plt.xlabel(}\StringTok{"Date"}\NormalTok{)}\OperatorTok{;}\NormalTok{ plt.ylabel(}\StringTok{"Valeur cumulée (base 1)"}\NormalTok{)}
\NormalTok{plt.legend()}\OperatorTok{;}\NormalTok{ plt.savefig(DOSSIER\_FIG }\OperatorTok{/} \StringTok{"04\_backtest.png"}\NormalTok{, dpi}\OperatorTok{=}\DecValTok{150}\NormalTok{, bbox\_inches}\OperatorTok{=}\StringTok{"tight"}\NormalTok{)}\OperatorTok{;}\NormalTok{ plt.show()}
\end{Highlighting}
\end{Shaded}

\begin{Shaded}
\begin{Highlighting}[]
\NormalTok{            stratégie  rendement cumulé (\%)  Sharpe annualisé  nb de changements}
\NormalTok{Régression logistique                 21.53              1.73                100}
\NormalTok{      Forêt aléatoire                 16.68              1.18                118}
\NormalTok{    Gradient boosting                 15.61              1.34                124}
\NormalTok{    Acheter{-}conserver                 32.28              1.68                  0}
\end{Highlighting}
\end{Shaded}

\begin{figure}
\centering
\includegraphics[width=1\textwidth,height=\textheight]{../../figures/04_backtest.png}
\caption{Backtest simplifié sur le test (coût 0,10 \%)}
\end{figure}

\textbf{Lecture.} Après coûts supposés de 0,10 \%, aucune stratégie ne
dépasse l'achat-conservation en rendement cumulé (+32,3 \%) ; la
régression logistique obtient +21,5 \% avec un Sharpe de 1,73 contre
1,68, écart trop faible pour être conclu sur 290 séances. Les stratégies
restent investies en moyenne moins longtemps dans un marché très
haussier, d'où leur retard.

\hypertarget{analyse-critique}{%
\subsection{8. Analyse critique}\label{analyse-critique}}

\begin{itemize}
\tightlist
\item
  \textbf{Réponse à la question.} La marche aléatoire \emph{stricte} est
  rejetée statistiquement pour le MASI (autocorrélation de rendement
  significative, ρ₁ ≈ 0,13), mais \textbf{nous ne démontrons pas qu'on
  puisse la battre de façon exploitable} : sur le test, aucun modèle ne
  dépasse la référence majoritaire ; en validation glissante, la
  meilleure (forêt, +2,8 points) n'est pas significative (p = 0,14).
\item
  \textbf{Puissance limitée.} Avec \textasciitilde290 observations,
  l'intervalle à 95 \% de la précision est d'environ ±5,6 points : un
  avantage réel de 2-3 points est indétectable. Absence de preuve ≠
  preuve d'absence.
\item
  \textbf{Sensibilité à la période.} La référence dépend de la part de
  hausses (50,7 \% train, 55,3 \% validation, 58,8 \% test) : un marché
  haussier favorise « toujours hausse ». Les résultats peuvent
  s'inverser dans un marché baissier.
\item
  \textbf{Comparaisons multiples.} Trois familles de modèles et
  \textasciitilde20 configurations ont été essayées : le risque de faux
  positif augmente ; les p-valeurs ne sont pas corrigées.
\item
  \textbf{Fuite résiduelle.} Les hyperparamètres choisis sur la
  validation sont réutilisés dans le \emph{walk-forward} (qui recouvre
  cette période) : légère optimisme possible. L'embargo d'une ligne
  traite le chevauchement de cible, mais pas la dépendance sérielle de
  la volatilité.
\item
  \textbf{Données.} Pas de volume, ouverture quasi toujours égale à la
  clôture de la veille, une seule série, indice non directement
  négociable (ETF/frais, liquidité, écarts de prix non modélisés). Le
  MASI peut subir un effet de cotation peu fréquente sur certaines
  valeurs, hypothèse plausible pour l'autocorrélation mais \textbf{non
  testée ici}.
\item
  \textbf{Coûts.} Les 0,10 \% sont une hypothèse ; la fiscalité et les
  frais réels peuvent dégrader davantage les stratégies actives (100-124
  changements de position sur le test).
\item
  \textbf{Pistes.} Classes de rendement avec seuil (ignorer les petites
  variations), modèles de volatilité (GARCH), plus d'actifs, correction
  des tests multiples.
\end{itemize}

\hypertarget{conclusion}{%
\subsection{9. Conclusion}\label{conclusion}}

Sur le MASI (2018-2025), \textbf{les modèles d'apprentissage supervisé
testés (régression logistique, forêt aléatoire, gradient boosting) ne
battent pas de façon significative la référence « marche aléatoire »}
(classe majoritaire) pour prédire le sens du rendement du lendemain :
sur le test, 50,3 \% à 53,4 \% contre 58,6 \% pour la référence ; en
walk-forward, au mieux 56,1 \% contre 53,3 \%, sans significativité (p =
0,14). L'objectif mesurable fixé en section 1 \textbf{n'est donc pas
atteint}. Le MASI présente néanmoins une autocorrélation de court terme
significative, qui justifie des études complémentaires plus puissantes
(plus d'actifs, plus de période, coûts réalistes).

\hypertarget{ruxe9fuxe9rences}{%
\subsection{10. Références}\label{ruxe9fuxe9rences}}

\begin{itemize}
\tightlist
\item
  Fama, E. F. (1970). \emph{Efficient Capital Markets: A Review of
  Theory and Empirical Work}. The Journal of Finance, 25(2), 383-417.
\item
  Lo, A. W. \& MacKinlay, A. C. (1988). \emph{Stock Market Prices Do Not
  Follow Random Walks: Evidence from a Simple Specification Test}. The
  Review of Financial Studies, 1(1), 41-66.
\item
  Ljung, G. M. \& Box, G. E. P. (1978). \emph{On a Measure of Lack of
  Fit in Time Series Models}. Biometrika, 65(2), 297-303.
\item
  McNemar, Q. (1947). \emph{Note on the sampling error of the difference
  between correlated proportions or percentages}. Psychometrika, 12(2),
  153-157.
\item
  Bergmeir, C. \& Benítez, J. M. (2012). \emph{On the use of
  cross-validation for time series predictor evaluation}. Information
  Sciences, 191, 192-213.
\item
  Pedregosa, F. et al.~(2011). \emph{Scikit-learn: Machine Learning in
  Python}. Journal of Machine Learning Research, 12, 2825-2830.
\item
  Hastie, T., Tibshirani, R. \& Friedman, J. (2009). \emph{The Elements
  of Statistical Learning} (2e éd.). Springer.
\item
  Données : fichier
  \texttt{Moroccan\_All\_Shares\_Historical\_Data\_\_1\_.csv} fourni
  pour le projet (indice MASI, Bourse de Casablanca ; export de type
  Investing.com).
\end{itemize}

\hypertarget{duxe9claration-des-ressources-et-outils-dia}{%
\subsection{Déclaration des ressources et outils
d'IA}\label{duxe9claration-des-ressources-et-outils-dia}}

\begin{quote}
\textbf{À relire et adapter par l'étudiant avant remise.}
\end{quote}

Un assistant d'IA générative (Claude, Anthropic) a été utilisé pour
produire une première version de la structure du notebook, du code
Python et des commentaires, à partir de la consigne et du fichier de
données fournis. Le code a été exécuté de bout en bout (Python 3.12,
graine 42) et les résultats présentés proviennent de cette exécution.
Les choix de conception (cible, variables, protocole chronologique,
tests) doivent être compris, vérifiés et assumés par l'auteur. Aucune
donnée externe n'a été utilisée en dehors du fichier fourni.

\end{document}
